\documentclass[11pt, a4paper]{report}

% Encoding & lingua
\usepackage[utf8]{inputenc}
\usepackage[T1]{fontenc}
\usepackage[italian]{babel}
\usepackage{underscore} % handles underscores well

% Impaginazione
\usepackage[
  top=2.5cm, bottom=2.5cm,
  left=3cm,  right=3cm
]{geometry}
\usepackage{microtype}
\usepackage{setspace}
\onehalfspacing

% Matematica
\usepackage{amsmath, amssymb, amsthm}
\usepackage{mathtools}
\usepackage[shortlabels]{enumitem}
\usepackage{lmodern}
\usepackage{bm}
\usepackage{tcolorbox}
% Ambienti teorema
\theoremstyle{plain}
\newtheorem{theorem}{Teorema}[chapter]
\newtheorem{proposition}[theorem]{Proposizione}
\newtheorem{lemma}[theorem]{Lemma}
\newtheorem{corollary}[theorem]{Corollario}
\theoremstyle{definition}
\newtheorem{definition}[theorem]{Definizione}
\newtheorem{esempio}[theorem]{Esempio}
\theoremstyle{remark}
\newtheorem{remark}[theorem]{Osservazione}

\usepackage{mdframed}
\usepackage{xcolor}

\newenvironment{hint}[1][Hint]
{\begin{mdframed}[backgroundcolor=blue!5, linecolor=blue!50, linewidth=2pt, topline=false, rightline=false, bottomline=false]
\textbf{\textsf{#1.}}\ }
{\end{mdframed}}

\newenvironment{spiegazione}[1][Spiegazione]
{\begin{mdframed}[backgroundcolor=green!5, linecolor=green!50, linewidth=2pt, topline=false, rightline=false, bottomline=false]
\textbf{\textsf{#1.}}\ }
{\end{mdframed}}

\newcommand{\PP}{\mathbb{P}}
\newcommand{\E}{\mathbb{E}}
\newcommand{\F}{\mathcal{F}}
\newcommand{\R}{\mathbb{R}}
\newcommand{\N}{\mathbb{N}}
\newcommand{\QQ}{\mathbb{Q}}
\newcommand{\1}{\mathbf{1}}

\begin{document}

\title{\Huge\textbf{Modelli Probabilistici}\\
\LARGE{Note per l'esame orale (Catene di Markov e Modello SIR)}}
\author{Francesco Castaldi}
\date{A.A. 2025/2026}

\maketitle

\tableofcontents
\newpage

\chapter{Teoria della probabilità in spazi finiti}

La teoria matematica della probabilità mira a costruire un formalismo teorico per descrivere i \emph{fenomeni aleatori} (casuali). Si basa sul concetto di \emph{esperimento aleatorio}, che può essere approssimativamente definito come "un esperimento il cui risultato non può essere previsto in anticipo". Un individuo può essere naturalmente portato ad associare un \emph{grado di fiducia} ai diversi possibili risultati di questo esperimento. Vediamo come modellizzare matematicamente questa idea, tenendo sempre presente che in ogni buona teoria vi sono tre fasi distinte:
\begin{enumerate}[(i)]
    \item Scelta del modello matematico;
    \item Elaborazione matematica all'interno del modello scelto;
    \item Interpretazione dei risultati ottenuti.
\end{enumerate}

\section{Spazi di probabilità}

La struttura degli spazi di probabilità si compone di tre ingredienti fondamentali:
\begin{enumerate}
    \item Lo spazio campionario (o spazio degli stati) $\Omega$;
    \item L'insieme degli eventi $\F$;
    \item La misura di probabilità $\PP$.
\end{enumerate}
Di conseguenza, uno spazio di probabilità è una terna: $(\Omega, \F, \PP)$. Descriviamo nel dettaglio gli oggetti di questa terna.

\begin{hint}[La Terna Fondamentale]
    Per ricordare la terna $(\Omega, \F, \PP)$:
    \begin{itemize}
        \item $\Omega$ (Omega) è \textbf{TUTTO} quello che può succedere (lo Spazio).
        \item $\F$ (F corsiva) sono le \textbf{DOMANDE} che puoi fare (gli Eventi a cui puoi assegnare una probabilità).
        \item $\PP$ (P) è la \textbf{RISPOSTA} (la Misura, il numero da 0 a 1).
    \end{itemize}
\end{hint}

\subsection{Spazio campionario ($\Omega$)}

Lo \emph{spazio campionario} è l'insieme i cui elementi rappresentano tutti i possibili "stati futuri/sconosciuti del mondo" \emph{mutuamente esclusivi}. È tradizionalmente indicato con $\Omega$ e il suo elemento generico è indicato con $\omega$. In questo capitolo assumiamo che lo spazio campionario $\Omega$ sia finito, ovvero $|\Omega| < \infty$. Naturalmente questa è un'assunzione restrittiva, ma ci permette comunque di modellare molti fenomeni.

\begin{esempio}[Monete e Dadi]
    \begin{itemize}
        \item \textbf{Lancio di una moneta:} Una scelta ragionevole per lo spazio campionario è $\Omega = \{T, C\}$ (Testa, Croce) oppure $\Omega = \{0, 1\}$. Le etichette non sono importanti.
        \item \textbf{Due lanci di una moneta (o lancio di due monete):} Una scelta logica è elencare le coppie: $\Omega := \{(T, T), (T, C), (C, T), (C, C)\}$.
        \item \textbf{Lancio di due dadi a sei facce:} $\Omega := \{(i, j) : i, j \in \{1, 2, 3, 4, 5, 6\}\}$. Si noti che $|\Omega| = 36$.
    \end{itemize}
\end{esempio}

\subsection{Spazio degli eventi ($\F$)}

Matematicamente, un \emph{evento} è un sottoinsieme $A \subseteq \Omega$ di risultati $\omega$. Interpretiamo $A$ come un evento che \emph{si è verificato} se il risultato effettivo $\omega \in A$, e che \emph{non si è verificato} se $\omega \notin A$. Questa rappresentazione matematica ci permette di tradurre domande intuitive in operazioni insiemistiche precise:
\begin{itemize}
    \item L'evento $\Omega$ è l'evento \emph{certo}.
    \item L'evento $\emptyset$ è l'evento \emph{impossibile}.
    \item Il singoletto $\{\omega\}$ è un evento \emph{elementare}.
    \item Dato $A \subseteq \Omega$, $A^c$ è l'"evento complementare (o contrario) di $A$".
    \item L'unione $A \cup B$ è l'evento "$A$ e/o $B$ si verificano".
    \item L'intersezione $A \cap B$ è l'evento "Sia $A$ che $B$ si verificano contemporaneamente".
\end{itemize}

L'insieme di tutti gli eventi di interesse è solitamente denotato con $\F$. Questo insieme deve avere delle "proprietà di stabilità", ovvero deve essere un'\emph{algebra}.

\begin{definition}[Algebra di insiemi e atomi]
Un insieme di sottoinsiemi $\F \subseteq 2^\Omega$ è chiamato \emph{algebra (di insiemi)} se valgono le seguenti proprietà:
\begin{enumerate}[(F1)]
    \item $\Omega \in \F$;
    \item $A \in \F \implies A^c \in \F$;
    \item $A, B \in \F \implies A \cup B \in \F$.
\end{enumerate}
Se $\F$ è un'algebra, un insieme $A \in \F$ è chiamato \emph{atomo} di $\F$ se non può essere decomposto come unione di due insiemi non vuoti e disgiunti di $\F$.
\end{definition}

\begin{spiegazione}
    Perché ci serve un'algebra? Perché se possiamo chiederci qual è la probabilità che piova (evento $A$), dobbiamo poterci chiedere anche qual è la probabilità che NON piova (evento $A^c$), oppure che piova o nevichi ($A \cup B$). L'algebra garantisce che il nostro "vocabolario" di eventi sia logicamente completo.
\end{spiegazione}

\subsection{Misura di probabilità ($\PP$)}

Una volta fissato lo spazio degli eventi $\F$, ad ogni evento $A \in \F$ si può associare un "grado di fiducia" chiamato \emph{probabilità dell'evento A}. Questo è formalizzato da una funzione:
\[ \PP : \F \to [0, 1] \]
La scelta di questa misura è soggettiva, ma deve rispettare i "tre assiomi della probabilità" per escludere scelte illogiche.

\begin{definition}[Misura di probabilità]
Sia $\Omega$ un insieme finito e sia $\F \subseteq 2^\Omega$ un'algebra. Una \emph{misura di probabilità} su $\F$ è una mappa $\PP : \F \to \R$ tale che:
\begin{enumerate}[(A1)]
    \item $0 \le \PP(A) \le 1$ per ogni $A \in \F$;
    \item $\PP(\Omega) = 1$;
    \item $A, B \in \F$ con $A \cap B = \emptyset \implies \PP(A \cup B) = \PP(A) + \PP(B)$.
\end{enumerate}
\end{definition}

\begin{hint}[Additività]
    L'assioma (A3) è la \textbf{Regola della Somma}: se due eventi non possono verificarsi insieme (sono disgiunti), la probabilità che se ne verifichi almeno uno è la somma delle loro probabilità singole. (Es. Probabilità di tirare 1 o 2 con un dado = P(1) + P(2)).
\end{hint}

Dagli assiomi derivano proprietà utilissime (prova a dimostrarle mentalmente usando i diagrammi di Eulero-Venn):
\begin{itemize}
    \item $\PP(A^c) = 1 - \PP(A)$
    \item Se $A \subseteq B$, allora $\PP(B \setminus A) = \PP(B) - \PP(A)$ e quindi $\PP(A) \le \PP(B)$.
    \item $\PP(A \cup B) = \PP(A) + \PP(B) - \PP(A \cap B)$
\end{itemize}

\section{Variabili aleatorie}

Nelle applicazioni pratiche, spesso non ci interessa l'esito esatto $\omega$, ma un numero o una quantità derivata da esso (es. la somma di due dadi, il guadagno in un gioco). Questo ci porta alle variabili aleatorie.

\begin{definition}[Variabile aleatoria]
Sia dato uno spazio $(\Omega, \F)$. Una mappa $X : \Omega \to E$ è una \emph{variabile aleatoria} se è $\F$-misurabile, ovvero se:
\[ X^{-1}(B) := \{X \in B\} \in \F \quad \forall B \subseteq E \]
\end{definition}

In spazi campionari finiti in cui $\F = 2^\Omega$ (l'insieme delle parti, ovvero tutti i possibili sottoinsiemi), \textbf{ogni funzione $X$ è una variabile aleatoria}.

\begin{spiegazione}
    Una variabile aleatoria non è né una "variabile" in senso stretto, né tantomeno "aleatoria" (casuale). È semplicemente una \textbf{funzione deterministica} che prende in input lo stato del mondo $\omega$ (quello sì che è casuale) e restituisce un valore, di solito un numero. \\
    Esempio: se $\omega = (\text{Testa}, \text{Testa})$, la variabile aleatoria $X = \text{"numero di teste"}$ restituirà banalmente $X(\omega) = 2$.
\end{spiegazione}

\subsection{Misurabilità di una mappa e algebra generata}

Il concetto di "algebra generata" serve a formalizzare l'idea di "informazione trasportata da una variabile aleatoria".

\begin{definition}[Algebra generata $\sigma(X)$]
Sia $X : \Omega \to E$ una variabile aleatoria. L'\emph{algebra generata} da $X$, denotata con $\sigma(X)$, è la più piccola algebra su $\Omega$ che rende $X$ misurabile:
\[ \sigma(X) := \{ X^{-1}(B) : B \subseteq E \} \]
\end{definition}

Gli atomi di $\sigma(X)$ sono esattamente gli insiemi di livello $\{X = x\}$ per ogni $x \in E$. 
\textbf{Intuizione:} L'algebra $\sigma(X)$ rappresenta tutta l'informazione che deduco osservando $X$.

\begin{theorem}[Criterio di misurabilità di Doob]
Siano $X : \Omega \to E$ e $Y : \Omega \to E'$ variabili aleatorie. 
$X$ è $\sigma(Y)$-misurabile se e solo se esiste una funzione $f : E' \to E$ tale che $X = f \circ Y$.
\end{theorem}

\begin{hint}[Interpretare Doob]
    Cosa significa che $X$ è $\sigma(Y)$-misurabile? 
    Significa che \textbf{l'informazione fornita da $Y$ è sufficiente a determinare $X$}. Se conosco $Y$, conosco automaticamente anche $X$. 
    Perché? Perché se $X = f(Y)$, basta applicare la funzione deterministica $f$ al valore di $Y$ per trovare $X$. Se invece serve lanciare una moneta aggiuntiva, allora $X$ \textbf{non} è $\sigma(Y)$-misurabile.
\end{hint}

\subsection{Legge (o distribuzione) di una variabile aleatoria}

\begin{definition}[Legge o distribuzione]
Data una variabile aleatoria $X : \Omega \to E$, possiamo trasferire la misura di probabilità da $(\Omega, \F, \PP)$ allo spazio dei valori $(E, 2^E)$ definendo:
\[ \mu_X(B) := \PP(X \in B) \quad \forall B \subseteq E \]
Questa nuova misura di probabilità si chiama \emph{legge} o \emph{distribuzione} di $X$.
La funzione $f_X(x) := \PP(X = x)$ è la sua \emph{funzione di densità discreta}.
\end{definition}

Poiché stiamo assumendo che $E$ sia finito, la distribuzione di probabilità (o legge) di una variabile aleatoria $X : \Omega \to E$ è completamente determinata dalla sua funzione di densità discreta $f_X$. Infatti, per ogni $B \subseteq E$, possiamo recuperare la legge $\mu_X(B)$ semplicemente sommando:
\[ \mu_X(B) = \sum_{x \in B} f_X(x) \]
Questa proprietà additiva rende le distribuzioni su spazi finiti particolarmente trattabili. Mentre la legge fornisce la misura di probabilità astratta, la funzione di densità discreta ci dà uno strumento computazionale concreto.

\begin{esempio}[Somma di due dadi]
Riprendendo il lancio di due dadi a sei facce, consideriamo la variabile aleatoria $S(\omega) = i + j$ (la somma dei due dadi). Il codominio è $E = \{2, \dots, 12\}$. La legge $\mu_S$ è descritta dalla densità discreta $f_S$ calcolata contando i casi favorevoli:
\[ f_S(2) = \frac{1}{36}, \quad f_S(3) = \frac{2}{36}, \quad \dots \quad f_S(7) = \frac{6}{36}, \dots \]
\end{esempio}

\subsection{Distribuzione congiunta e distribuzioni marginali}

Quando studiamo più variabili aleatorie contemporaneamente, abbiamo bisogno di capire non solo il loro comportamento individuale, ma anche come si relazionano tra loro.

\begin{definition}[Leggi congiunte e marginali]
Date $n$ variabili aleatorie $X_j : \Omega \to E_j$, con $j=1,\dots,n$, il vettore $X = (X_1, \dots, X_n)$ è una variabile aleatoria che assume valori nello spazio prodotto $E = E_1 \times \dots \times E_n$.
La \emph{legge congiunta} $\mu_X = L(X_1, \dots, X_n)$ è la misura di probabilità su $E$ definita da:
\[ \mu_X(A_1 \times \dots \times A_n) := \PP(X_1 \in A_1, \dots, X_n \in A_n) \]
Le \emph{leggi marginali} $\mu_{X_i} = L(X_i)$ si ottengono dalla legge congiunta:
\[ \mu_{X_i}(A_i) = \mu_X(E_1 \times \dots \times A_i \times \dots \times E_n) \]
La corrispondente \emph{densità discreta marginale} di $X_i$ si ottiene sommando rispetto a tutte le altre variabili:
\[ f_{X_i}(x_i) = \sum_{x_1, \dots, x_{i-1}, x_{i+1}, \dots, x_n} f_X(x_1, \dots, x_n) \]
\end{definition}

\begin{spiegazione}
    Immagina la legge congiunta come una grande tabella a doppia entrata (se abbiamo due variabili $X$ e $Y$). Le celle interne della tabella contengono la probabilità che $X=x$ e $Y=y$ contemporaneamente. 
    Se sommiamo tutti i valori di una riga, troviamo la probabilità totale di $X=x$ (a prescindere da $Y$). Questa somma "ai margini" della tabella è proprio la \emph{distribuzione marginale}.
\end{spiegazione}

\begin{hint}[Congiunta vs Marginale]
    Ricorda: \textbf{Dalla congiunta puoi SEMPRE ricavare le marginali (sommando), ma dalle marginali NON puoi ricavare la congiunta} (a meno che le variabili non siano indipendenti). Diverse distribuzioni congiunte possono avere le stesse identiche distribuzioni marginali.
\end{hint}

\section{Probabilità Condizionata}

Esploriamo ora due concetti fondamentali che catturano come le probabilità cambiano quando arrivano nuove informazioni.

\begin{definition}[Probabilità Condizionata]
Dato uno spazio di probabilità $(\Omega, \F, \PP)$ e un evento $H \in \F$ con $\PP(H) > 0$, la \emph{probabilità condizionata} di $A$ dato $H$ è definita come:
\[ \PP(A|H) := \frac{\PP(A \cap H)}{\PP(H)} \]
\end{definition}

La probabilità condizionata rappresenta l'aggiornamento razionale delle credenze di un agente quando viene a sapere che l'evento $H$ si è verificato. L'operazione "rinormalizza" la probabilità originale concentrandola su $H$.

\begin{esempio}[Il problema di Monty Hall]
Sei in un gioco a premi e devi scegliere tra 3 porte: dietro una c'è un'auto, dietro le altre due ci sono capre. Scegli la porta 1. Il presentatore (che sa cosa c'è dietro le porte) apre la porta 3, mostrando una capra. Ti chiede: "Vuoi cambiare e scegliere la porta 2?".
\textbf{Conviene cambiare?}
Sì, le probabilità raddoppiano (da 1/3 a 2/3). Il motivo è che l'azione del presentatore non è casuale: è costretto ad aprire una porta con una capra. Se la tua prima scelta era sbagliata (probabilità 2/3), il presentatore ti sta indicando esattamente l'altra porta sbagliata, lasciando l'auto dietro l'unica porta chiusa rimasta. Cambiando vinci sempre in questo caso!
\end{esempio}

\subsection{Legge della probabilità totale e Teorema di Bayes}

Queste formule sono fondamentali per scomporre eventi complessi in parti più semplici o per invertire i condizionamenti.

\begin{theorem}[Probabilità totale e Bayes]
\begin{enumerate}
    \item \textbf{Legge della probabilità totale}: Se $A_1, \dots, A_n$ formano una partizione di $\Omega$, per ogni evento $A$:
    \[ \PP(A) = \sum_{i=1}^n \PP(A_i) \PP(A|A_i) \]
    \item \textbf{Teorema di Bayes}: Per due eventi $A$ e $B$ con $\PP(A)>0$ e $\PP(B)>0$:
    \[ \PP(A|B) = \frac{\PP(B|A)\PP(A)}{\PP(B)} \]
\end{enumerate}
\end{theorem}

\begin{spiegazione}
    \textbf{Probabilità Totale}: Se voglio calcolare la probabilità che piova (evento A), posso dividere il mondo in due scenari (partizione): $A_1$=è nuvoloso, $A_2$=è sereno. La probabilità che piova sarà la probabilità che piova QUANDO è nuvoloso (moltiplicata per la probabilità che sia nuvoloso) PIÙ la probabilità che piova QUANDO è sereno (moltiplicata per la probabilità che sia sereno). \\
    \textbf{Bayes}: Serve a "ribaltare" la domanda. Conosciamo la probabilità di avere un test positivo sapendo di essere malati $P(+|Malato)$, ma a noi interessa la probabilità di essere malati avendo il test positivo $P(Malato|+)$! Bayes fa esattamente questo.
\end{spiegazione}

\begin{hint}[Formula di Bayes in Pratica]
    Quando applichi Bayes, spesso il denominatore $\PP(B)$ non è dato direttamente. Ricorda di espanderlo sempre usando la Legge della Probabilità Totale al denominatore: 
    \[ \PP(A_1|B) = \frac{\PP(B|A_1)\PP(A_1)}{\PP(B|A_1)\PP(A_1) + \PP(B|A_2)\PP(A_2) + \dots} \]
\end{hint}

\section{Indipendenza}

L'indipendenza cattura le situazioni in cui la conoscenza di un fenomeno casuale non fornisce alcuna informazione su un altro. Intuitivamente, non c'è "interazione".

\begin{definition}[Indipendenza tra variabili aleatorie]
Due variabili aleatorie $X : \Omega \to E$ e $Y : \Omega \to E'$ sono \emph{indipendenti} se per ogni evento della forma $H = \{Y \in B\}$ con $B \subseteq E'$ e $\PP(H)>0$, si ha:
\[ \PP(X \in A | Y \in B) = \PP(X \in A) \]
In modo equivalente, c'è una fattorizzazione simmetrica:
\[ \PP(X \in A, Y \in B) = \PP(X \in A)\PP(Y \in B) \]
\end{definition}

\begin{definition}[Indipendenza tra eventi]
Due eventi $A, B \in \F$ sono indipendenti se e solo se:
\[ \PP(A \cap B) = \PP(A)\PP(B) \]
Questo equivale a dire che le loro funzioni indicatrici $\mathbf{1}_A$ e $\mathbf{1}_B$ sono variabili aleatorie indipendenti.
\end{definition}
\begin{theorem}[Caratterizzazione dell'indipendenza]
Due variabili aleatorie $X$ e $Y$ sono indipendenti se e solo se la loro distribuzione congiunta è il prodotto delle marginali:
\[ \mu_{(X,Y)}(A \times B) = \mu_X(A)\mu_Y(B) \]
Oppure, in termini di densità:
\[ f_{X,Y}(x, y) = f_X(x)f_Y(y) \]
\end{theorem}

\subsection{Famiglie indipendenti}

L'indipendenza a due a due non basta quando abbiamo molte variabili.

\begin{definition}[Indipendenza di una famiglia]
Una famiglia finita di variabili aleatorie $\{X_1, \dots, X_n\}$ è \emph{indipendente} se per qualsiasi due sottofamiglie disgiunte, i vettori risultanti sono indipendenti. Questo equivale a dire che la densità congiunta fattorizza in tutte le marginali:
\[ f_{X_1, \dots, X_n}(x_1, \dots, x_n) = f_{X_1}(x_1) \dots f_{X_n}(x_n) \]
\end{definition}

\begin{spiegazione}
    Un errore comune è pensare che se tutte le coppie di variabili in un gruppo sono indipendenti tra loro (indipendenza a due a due), allora tutto il gruppo è indipendente. L'esempio di Bernstein ci mostra che questo è falso! Due monete possono essere indipendenti, ma una terza moneta il cui risultato dipende dal prodotto delle prime due rompe l'indipendenza globale, pur mantenendola a coppie.
\end{spiegazione}

\begin{theorem}[Conservazione dell'indipendenza per trasformazioni]
Se $\{X_1, \dots, X_n\}$ sono variabili indipendenti, allora applicando delle funzioni deterministiche (es. elevarle al quadrato, moltiplicarle per costanti), le nuove variabili risultanti $\{f_1(X_1), \dots, f_n(X_n)\}$ \textbf{rimangono indipendenti}. L'elaborazione deterministica non "crea" correlazione dal nulla.
\end{theorem}

\section{Tre distribuzioni fondamentali (su insiemi finiti)}

\subsection{Distribuzione Uniforme}
Modella esperimenti in cui un oggetto è scelto in modo del tutto casuale ed equo da un insieme di $n$ elementi.
La densità discreta è:
\[ f_X(x) = \frac{1}{n} \]

\subsection{Distribuzione di Bernoulli}
Modella un singolo esperimento binario (es. lancio di moneta), con esito 1 (successo) o 0 (fallimento).
Indicata con $\mathcal{B}(p)$ dove $p \in (0,1)$ è la probabilità di successo.
\[ f_X(1) = p, \quad f_X(0) = 1-p \]

\subsection{Distribuzione Binomiale}
Conta il \textbf{numero totale di successi} su $n$ esperimenti di Bernoulli \emph{indipendenti} e con la stessa probabilità $p$.
Denotata con $\mathcal{B}(n, p)$.
Per calcolare la probabilità di avere esattamente $i$ successi, si usa il calcolo combinatorio (in particolare il coefficiente binomiale $\binom{n}{i}$):
\[ \PP(S = i) = \binom{n}{i} p^i (1-p)^{n-i} \]

\begin{hint}[Binomiale = Somma di Bernoulli]
    È fondamentale ricordare che una variabile Binomiale $S \sim \mathcal{B}(n,p)$ può SEMPRE essere scritta come somma di $n$ variabili di Bernoulli indipendenti: $S = X_1 + X_2 + \dots + X_n$. Questo "trucco" rende i calcoli (es. il valore atteso) facilissimi usando la linearità!
\end{hint}

\section{Valore Atteso e Varianza}

Il \emph{valore atteso} (o media) è il concetto fondamentale per quantificare la "tendenza" centrale di una variabile aleatoria.

\begin{definition}[Valore atteso]
Sia $X : \Omega \to \R$ una variabile aleatoria. Il suo valore atteso è:
\[ \E[X] := \sum_{x \in E} x \PP(X = x) = \sum_{x \in E} x f_X(x) \]
\end{definition}

\begin{spiegazione}
    Il valore atteso NON è il valore che ti "aspetti" di vedere in un singolo esperimento (può anche essere un numero impossibile da ottenere, come 3.5 per un dado). È la \textbf{media ponderata} di tutti i possibili valori, dove i pesi sono le probabilità. Rappresenta la migliore stima costante della variabile.
\end{spiegazione}

Le proprietà chiave del valore atteso sono la \textbf{linearità} e la \textbf{monotonia}:
\[ \E[\alpha X + \beta Y] = \alpha \E[X] + \beta \E[Y] \]

\begin{definition}[Varianza]
La varianza misura quanto una variabile si "discosta" dal suo valore medio. È definita come:
\[ \operatorname{Var}(X) := \E[(X - \E[X])^2] \]
\end{definition}

\section{Valore Atteso Condizionato}

Consideriamo ora il valore atteso di una variabile aleatoria rispetto a una probabilità condizionata.

\begin{definition}[Valore atteso condizionato a un evento]
Dato un evento $H \in \F$ con $\PP(H) > 0$, il \emph{valore atteso condizionato} di $X$ dato $H$ è:
\[ \E[X|H] := \sum_{x \in E} x \PP(X = x | H) \]
\end{definition}

Questo ci permette di definire una formula utilissima, speculare alla Legge della Probabilità Totale.

\begin{theorem}[Formula di disintegrazione]
Se $A_1, \dots, A_n$ è una partizione di $\Omega$, il valore atteso globale si può ricavare combinando i valori attesi condizionati:
\[ \E[X] = \sum_{i=1}^n \E[X|A_i] \PP(A_i) \]
\end{theorem}

\subsection{Valore atteso condizionato rispetto a una variabile aleatoria}

Possiamo anche condizionare rispetto a un'altra variabile aleatoria $Y$. In questo caso, il risultato non è un singolo numero, ma \textbf{una nuova variabile aleatoria} che dipende dai valori assunti da $Y$.

\begin{definition}
Il \emph{valore atteso condizionato} di $X$ data $Y$, denotato con $\E[X|Y]$, è la variabile aleatoria definita come:
\[ \E[X|Y](\omega) = \E[X | Y = Y(\omega)] \]
\end{definition}

\begin{spiegazione}
    Mentre $\E[X|Y=y]$ è un numero fisso (perché abbiamo "fissato" l'osservazione $Y=y$), $\E[X|Y]$ è una funzione deterministica di $Y$. Rappresenta la nostra "migliore stima" di $X$ in base a quale valore assume $Y$ nel mondo.
\end{spiegazione}

Proprietà fondamentali di $\E[X|Y]$:
\begin{itemize}
    \item Se $X$ e $Y$ sono indipendenti: la conoscenza di $Y$ non aiuta a stimare $X$, quindi $\E[X|Y] = \E[X]$.
    \item Se $X$ è completamente determinata da $Y$ (cioè $X = f(Y)$): la stima migliore è il valore esatto, quindi $\E[X|Y] = X$.
    \item \textbf{Legge delle aspettative iterate (Tower property)}: $\E[\E[X|Y]] = \E[X]$. In media, la nostra stima condizionata ci restituisce il valore atteso globale!
\end{itemize}

\begin{theorem}[Freezing lemma]
Sia $\phi(X,Y)$ una funzione di due variabili. Il valore atteso condizionato di $\phi(X,Y)$ data $Y$ si ottiene "congelando" $Y$ al suo valore osservato e calcolando la media solo rispetto a $X$:
\[ \E[\phi(X,Y)|Y](\omega) = \E[\phi(X, y)]\Big|_{y=Y(\omega)} \]
\end{theorem}


\chapter{Il caso discreto: da finito a numerabile}

Fino ad ora abbiamo considerato solo spazi di probabilità finiti. Questo è un limite significativo per descrivere molti fenomeni naturali. Ad esempio, non potremmo descrivere una variabile aleatoria che conta il numero di lanci necessari per ottenere testa, poiché i lanci potrebbero essere teoricamente infiniti, richiedendo un codominio come $\N$. Allo stesso modo, non potremmo descrivere processi su orizzonti temporali infiniti. 

Dobbiamo estendere il nostro quadro teorico dai casi finiti ai casi infiniti, in particolare spazi numerabili (o più grandi). Ecco le estensioni chiave che affronteremo:
\begin{itemize}
    \item \textbf{Spazio campionario} $\Omega$: da finito a infinito.
    \item \textbf{$\sigma$-algebra}: chiusura rispetto alle unioni da finita a \emph{numerabile}.
    \item \textbf{Misura di probabilità}: dall'additività finita all'\emph{additività numerabile} ($\sigma$-additività).
    \item \textbf{Famiglie di variabili aleatorie}: da collezioni finite a collezioni \emph{numerabili} $\{X_j\}_{j \in \N}$.
\end{itemize}

\section{Spazi di probabilità generali}

\begin{definition}[$\sigma$-algebra]
Una collezione $\F \subseteq 2^\Omega$ è chiamata \emph{$\sigma$-algebra} su $\Omega$ se:
\begin{enumerate}
    \item $\Omega \in \F$,
    \item $A \in \F \implies A^c \in \F$ (chiusura rispetto al complemento),
    \item Per ogni collezione numerabile $\{A_n\}_{n \in \N} \subseteq \F$, l'unione $\bigcup_{n \in \N} A_n \in \F$ (chiusura rispetto all'unione numerabile).
\end{enumerate}
La coppia $(\Omega, \F)$ è detta \emph{spazio misurabile}.
\end{definition}

\begin{definition}[Misura di probabilità ($\sigma$-additiva)]
Sia $\F$ una $\sigma$-algebra. Una funzione $\PP : \F \to [0, 1]$ è una \emph{misura di probabilità} su $(\Omega, \F)$ se:
\begin{enumerate}
    \item $0 \le \PP(A) \le 1$ per ogni $A \in \F$,
    \item $\PP(\Omega) = 1$,
    \item \textbf{$\sigma$-additività}: per ogni collezione numerabile di eventi \emph{disgiunti} a due a due $\{A_n\}_{n \in \N}$, vale:
    \[ \PP\left(\bigcup_{n=1}^\infty A_n\right) = \sum_{n=1}^\infty \PP(A_n) \]
\end{enumerate}
\end{definition}

\begin{spiegazione}
    Il passaggio dall'additività finita alla $\sigma$-additività (numerabile) è il cuore della probabilità moderna introdotta da Kolmogorov. Ci permette di sommare probabilità di infiniti eventi disgiunti e garantisce che la probabilità si comporti bene quando facciamo limiti di eventi (continuità dal basso e dall'alto).
\end{spiegazione}

\section{Variabili aleatorie discrete}

Limitiamo la nostra teoria alle variabili aleatorie che hanno un codominio \emph{al più numerabile}. Queste sono chiamate \textbf{variabili aleatorie discrete}. Tutti i concetti di densità discreta e distribuzione del Capitolo 1 si estendono naturalmente sostituendo le somme finite con le \emph{serie} numerabili.

\subsection{Distribuzione di Poisson}

Una distribuzione fondamentale sugli interi naturali $\N$ è la distribuzione di Poisson.

\begin{definition}[Distribuzione di Poisson]
La distribuzione di Poisson di parametro $\lambda > 0$, indicata con $\mathcal{P}(\lambda)$, è la legge di probabilità su $\N$ definita da:
\[ \PP(X = k) = e^{-\lambda} \frac{\lambda^k}{k!}, \quad k \in \N \]
\end{definition}

\begin{hint}[Somma di Poisson]
    Ricorda che se sommiamo due variabili di Poisson indipendenti con parametri $\lambda_1$ e $\lambda_2$, il risultato è ancora una variabile di Poisson con parametro $\lambda_1 + \lambda_2$.
\end{hint}

\section{Schema di Bernoulli e sequenze indipendenti}

Vogliamo modellare un esperimento che consiste in infiniti lanci indipendenti di una moneta. Abbiamo bisogno della garanzia matematica che esista uno spazio di probabilità abbastanza grande da supportare una sequenza infinita di variabili indipendenti.

\begin{theorem}[Teorema di estensione di Kolmogorov - Schema di Bernoulli]
Sia $(E, \mathcal{E})$ uno spazio discreto e $\mu$ una distribuzione di probabilità. Esiste uno spazio di probabilità $(\Omega, \F, \PP)$ e una sequenza \textbf{infinita} di variabili aleatorie \textbf{indipendenti e identicamente distribuite (i.i.d.)} $X_1, X_2, \dots$ a valori in $E$ tale che ogni $X_i$ abbia distribuzione $\mu$.
\end{theorem}

Questo teorema legittima l'uso di infinite prove ripetute (lo schema di Bernoulli).

\subsection{Distribuzione Geometrica}

Consideriamo una sequenza infinita di tentativi di Bernoulli i.i.d.\ con probabilità di successo $p \in (0,1)$. Sia $T$ la variabile aleatoria che conta il \emph{numero di tentativi fino al primo successo} (incluso). 

\begin{definition}[Distribuzione Geometrica]
La legge del "tempo di primo successo" $T$ si chiama \emph{distribuzione geometrica di parametro $p$}. La sua densità discreta è:
\[ \PP(T = n) = (1-p)^{n-1} p, \quad n \in \N, n \ge 1 \]
\end{definition}
(Notiamo che la probabilità che non ci sia mai un successo è zero, $\PP(T=\infty)=0$).

\begin{theorem}[Proprietà di Assenza di Memoria]
La distribuzione geometrica è l'unica distribuzione discreta ad avere la proprietà di assenza di memoria (\emph{memoryless property}):
\[ \PP(T = n + k | T > n) = \PP(T = k) \quad \forall n, k \ge 1 \]
\end{theorem}

\begin{spiegazione}
    L'assenza di memoria significa che "la moneta non ha memoria". Se hai appena fatto 10 lanci fallimentari e aspetti il successo, la probabilità di dover aspettare altri 3 lanci è IDENTICA alla probabilità che avresti avuto all'inizio di dover aspettare 3 lanci. Il passato non cambia l'attesa futura!
\end{spiegazione}

\section{Valore Atteso nel Caso Discreto Generale}

Quando passiamo a spazi numerabili, le somme diventano serie infinite. Dobbiamo fare attenzione alla convergenza. 

\begin{definition}[Variabile aleatoria semi-integrabile e valore atteso]
Sia $X : \Omega \to \R$ una variabile aleatoria discreta (a valori in un insieme numerabile). Definiamo il suo \emph{valore atteso} separando la parte positiva $X^+$ dalla parte negativa $X^-$:
\[ \E[X] = \E[X^+] - \E[X^-] \]
Se almeno uno dei due termini è finito, diciamo che $X$ è \emph{semi-integrabile}. Se entrambi sono finiti, $X$ è \emph{integrabile} e scriviamo $X \in L^1(\Omega, \F, \PP)$.
\end{definition}

Le proprietà fondamentali del valore atteso (Linearità, Monotonia, Positività, Legge dello Statistico Inconsapevole) rimangono identiche, a patto di garantire la convergenza assoluta delle serie coinvolte (cioè $X, Y \in L^1$).

Un risultato estremamente utile per variabili a valori interi positivi è il seguente:
\begin{theorem}[Formula alternativa per il valore atteso intero]
Se $X$ assume valori in $\N \cup \{0, \infty\}$, il suo valore atteso può essere calcolato usando le probabilità cumulate complementari (code della distribuzione):
\[ \E[X] = \sum_{n=1}^\infty \PP(X \ge n) \]
\end{theorem}

\subsection{Varianza e Disuguaglianze Notevoli}
La Varianza si definisce come $\operatorname{Var}(X) = \E[(X-\E[X])^2]$ per variabili con varianza finita. In questo contesto, le disuguaglianze di Markov e Chebyshev sono strumenti fondamentali per stimare probabilità anche senza conoscere la densità esatta.

\begin{theorem}[Disuguaglianza di Markov e di Chebyshev]
\textbf{Markov:} Per ogni $X$ discreta non negativa e per ogni $a > 0$:
\[ \PP(X \ge a) \le \frac{\E[X]}{a} \]

\textbf{Chebyshev:} Per ogni $X$ discreta con media $m$ e varianza $\sigma^2$ finita, e per ogni $\epsilon > 0$:
\[ \PP(|X - m| \ge \epsilon) \le \frac{\operatorname{Var}(X)}{\epsilon^2} \]
\end{theorem}

\section{Valore Atteso Condizionato (Caso Generale)}

Anche il valore atteso condizionato $\E[X|Y]$ (inteso come variabile aleatoria) mantiene le sue proprietà fondamentali, con l'accortezza di sostituire le somme finite con le serie.
Ricordiamo le proprietà più importanti, dove si assume che tutte le variabili siano in $L^1$:
\begin{enumerate}
    \item \textbf{Legge delle aspettative iterate (Tower Property):} $\E[\E[X|Y]] = \E[X]$.
    \item \textbf{Tirar fuori ciò che è noto (Taking out what is known):} $\E[X Z | Y] = Z \E[X | Y]$ se $Z$ è una funzione deterministica di $Y$.
    \item \textbf{Disuguaglianza di Jensen Condizionata:} Se $\phi$ è convessa e $X \in L^1$, allora $\phi(\E[X|Y]) \le \E[\phi(X)|Y]$.
    \item \textbf{Freezing Lemma:} $\E[\phi(X,Y)|Y](\omega) = \E[\phi(X, y)]\big|_{y=Y(\omega)}$.
\end{enumerate}

\begin{hint}[La formula di disintegrazione e il Freezing Lemma]
    Combinando il Freezing Lemma con la Tower Property otteniamo una generalizzazione della formula di disintegrazione:
    \[ \E[\phi(X,Y)] = \E[\E[\phi(X,Y)|Y]] \]
    Se $X$ e $Y$ sono indipendenti, calcolare la media interna è facilissimo perché $\E[\phi(X,y)|Y=y] = \E[\phi(X,y)]$.
\end{hint}

\begin{esempio}[Equazione di Wald]
Sia $(X_i)_{i \ge 1}$ una successione di variabili i.i.d. con media $\mu$, e sia $N$ una variabile a valori in $\N$, \emph{indipendente} dalle $X_i$. Definiamo la somma aleatoria $S_N = \sum_{i=1}^N X_i$.
Qual è il suo valore atteso? Usando il condizionamento rispetto a $N$:
\[ \E[S_N|N=n] = \E\left[\sum_{i=1}^n X_i\right] = n\mu \]
Quindi $\E[S_N|N] = N\mu$. Applicando la Tower Property: $\E[S_N] = \E[N\mu] = \mu \E[N]$. Questa è nota come Equazione di Wald.
\end{esempio}

\section{Convergenze di Variabili Aleatorie}

Una delle necessità più grandi nel passare a collezioni numerabili di variabili (es. le medie di $n$ esperimenti per $n \to \infty$) è poter studiare il limite delle variabili aleatorie.
In probabilità, abbiamo tre concetti principali di convergenza.

\begin{definition}[Tipi di Convergenza]
Data una successione $X_n$ e una variabile limite $X$:
\begin{enumerate}
    \item \textbf{Convergenza Quasi Certamente (a.s.):} $X_n \to X$ q.c. se $\PP(\{\omega \in \Omega : \lim_{n\to\infty} X_n(\omega) = X(\omega)\}) = 1$. (Per quasi ogni esito, la sequenza numerica converge).
    \item \textbf{Convergenza in Probabilità:} $X_n \stackrel{P}{\to} X$ se per ogni $\epsilon > 0$, $\lim_{n\to\infty} \PP(|X_n - X| \ge \epsilon) = 0$. (La probabilità che $X_n$ disti da $X$ più di $\epsilon$ svanisce).
    \item \textbf{Convergenza in Legge (o in Distribuzione):} $X_n \stackrel{\mathcal{L}}{\to} X$ se per ogni funzione $\phi$ limitata e continua, $\E[\phi(X_n)] \to \E[\phi(X)]$. (Le \emph{distribuzioni} di $X_n$ diventano simili a quella di $X$, ma le variabili stesse su $\Omega$ potrebbero non essere "vicine").
\end{enumerate}
\end{definition}

\begin{spiegazione}
    Esiste una gerarchia rigida tra queste convergenze:
    \[ \text{Quasi Certamente} \implies \text{In Probabilità} \implies \text{In Legge} \]
    Le implicazioni inverse \textbf{non} valgono in generale. 
    Ad esempio, se lancio all'infinito due dadi sempre diversi ma entrambi bilanciati, le loro distribuzioni (leggi) sono identiche, ma il risultato del dado $n$-esimo non "converge numericamente" in probabilità al risultato di un dado $X$.
\end{spiegazione}

\subsection{Teoremi di convergenza in media}

Passare al limite dentro un valore atteso ($\lim \E[X_n] = \E[\lim X_n]$) non è sempre possibile. Esistono tre grandi teoremi per legittimare questo passaggio, mutuati dalla teoria dell'integrazione di Lebesgue:

\begin{theorem}[Teorema di Convergenza Monotona (Beppo-Levi)]
Se $X_n \ge 0$ e $X_n \uparrow X$ monotonamente (quasi certamente), allora:
\[ \lim_{n\to\infty} \E[X_n] = \E[X] \]
\end{theorem}

\begin{theorem}[Lemma di Fatou]
Se $X_n \ge 0$ (quasi certamente) per ogni $n$, allora l'aspettativa del limite inferiore è minore o uguale al limite inferiore delle aspettative:
\[ \E\left[\liminf_{n\to\infty} X_n\right] \le \liminf_{n\to\infty} \E[X_n] \]
\end{theorem}

\begin{theorem}[Teorema di Convergenza Dominata (Lebesgue)]
Se $X_n \to X$ (in probabilità o quasi certamente) e se esiste una variabile "dominante" $G \in L^1$ tale che $|X_n| \le G$ per ogni $n$, allora:
\[ \lim_{n\to\infty} \E[X_n] = \E[X] \]
\end{theorem}

\section{La Legge dei Grandi Numeri}

Uno dei risultati più celebri della teoria della probabilità è la \emph{Legge dei Grandi Numeri}. In senso lato, afferma che:
\emph{"La media empirica di una successione di variabili aleatorie identicamente distribuite converge (in un senso opportuno) al loro valore atteso comune."}

Esistono due versioni principali, a seconda del tipo di convergenza (e delle ipotesi richieste).

\begin{theorem}[Legge debole dei grandi numeri]
Sia $(X_n)_{n \in \N}$ una successione di variabili i.i.d. discrete a valori reali tali che $\E[X_1^2] < \infty$ (il che implica varianza finita $\sigma^2$). Sia $m = \E[X_1]$. Posto:
\[ \bar{S}_n = \frac{X_1 + \cdots + X_n}{n} \]
allora per ogni $\epsilon > 0$:
\[ \lim_{n\to\infty} \PP(|\bar{S}_n - m| \ge \epsilon) = 0 \]
(cioè $\bar{S}_n \stackrel{P}{\to} m$).
\end{theorem}
\emph{Dimostrazione (Idea):} Usando la linearità, $\E[\bar{S}_n] = m$. Per l'indipendenza, $\operatorname{Var}(\bar{S}_n) = \frac{\sigma^2}{n}$. Applicando Chebyshev, $\PP(|\bar{S}_n - m| \ge \epsilon) \le \frac{\sigma^2}{n\epsilon^2}$, che tende a 0.

\begin{theorem}[Legge forte dei grandi numeri di Borel]
Sia $(X_n)_{n \in \N}$ i.i.d. e \textbf{limitate} (esiste $C$ tale che $|X_i| \le C$). Posto $m = \E[X_1]$, si ha:
\[ \PP\left( \lim_{n\to\infty} \frac{S_n}{n} = m \right) = 1 \]
(cioè $\bar{S}_n \to m$ quasi certamente).
\end{theorem}

\begin{hint}[Riflessione Frequenzista e Assiomatica]
È fondamentale evitare la circolarità logica frequenzista: \textbf{la probabilità NON è il limite delle frequenze}. Se lo fosse, la Legge dei Grandi Numeri sarebbe una tautologia (una definizione travestita da teorema).
Invece, nell'assiomatica di Kolmogorov, la probabilità è una misura primitiva assegnata \emph{a priori}. La Legge dei Grandi Numeri ci dice che, sotto queste regole, il modello è internamente coerente: i risultati pratici a lungo termine si allineano con la costruzione matematica.
\end{hint}

\section{Esercizi Selezionati (Capitolo 2)}

Di seguito alcuni esercizi essenziali per fissare i concetti su variabili discrete, condizionamento e modelli classici.

\subsection*{Esercizio 1: Tempi d'attesa e parità}
Lanciando una moneta equa, qual è la probabilità che il primo successo avvenga in un lancio dispari?
\begin{spiegazione}
Sia $T \sim \text{Geom}(1/2)$. La probabilità è data dalla serie:
\[ \PP(T \in \{1,3,5,\dots\}) = \sum_{k=0}^\infty \PP(T = 2k+1) = \sum_{k=0}^\infty (1/2)^{2k+1} \]
Questo è pari a $\frac{1}{2} \sum (1/4)^k = \frac{1/2}{1 - 1/4} = \frac{1/2}{3/4} = \frac{2}{3}$.
\end{spiegazione}

\subsection*{Esercizio 2: Estrazioni con stop al primo disaccordo}
Due giocatori (A e B) estraggono (con rimpiazzo) da urne identiche ($r$ rosse, $b$ nere). Il gioco si ferma al primo turno in cui i risultati divergono. Chi ha estratto rosso, vince 1 Euro. Qual è la legge della vincita $V$ di A, e del numero di turni $N$? Sono indipendenti?
\begin{spiegazione}
La probabilità di successo al turno $n$ per A è $p = r/(r+b)$. Il gioco si ferma se (A estrae rosso e B nera) o viceversa. La probabilità che fermi a un turno è $2p(1-p)$.
$N$ ha legge geometrica di parametro $2p(1-p)$.
Dato che il gioco si è fermato, la probabilità che a vincere sia stato A (rosso) o B (rosso) è per simmetria esattamente $1/2$.
Quindi $V \sim B(1/2)$ e il risultato finale è \textbf{indipendente} dalla durata $N$.
\end{spiegazione}

\subsection*{Esercizio 3: Tempi di successo successivi}
In uno schema di Bernoulli con probabilità $p$, se so che il secondo successo è avvenuto al tempo $N$ ($T_2 = N$), qual è la legge del primo successo $T_1$?
\begin{spiegazione}
Usiamo Bayes: $\PP(T_1 = k | T_2 = N) \propto \PP(T_1 = k, T_2 = N)$.
L'evento $\{T_1 = k, T_2 = N\}$ equivale a: "1 successo al tempo $k$, insuccessi ovunque altro prima di $N$, e successo al tempo $N$". La probabilità è $p^2 (1-p)^{N-2}$, che \textbf{non dipende da $k$}.
Quindi, dato $T_2=N$, il tempo $T_1$ è uniformemente distribuito su $\{1, \dots, N-1\}$.
\end{spiegazione}

\subsection*{Esercizio 4: Estrazioni senza rimpiazzo (Variabili ipergeometriche)}
Estraggo $k$ palline su $n$ totali senza rimpiazzo. Qual è la probabilità di estrarre un elemento specifico $x$? E due elementi $x, y$?
\begin{spiegazione}
Sia $X_i$ l'esito della $i$-esima estrazione. Per simmetria, ogni estrazione marginalmente è uniforme: $\PP(X_i = x) = 1/n$.
La probabilità che $x$ sia estratto è $\PP(\cup_{i=1}^k \{X_i=x\}) = \sum \PP(X_i=x) = k/n$.
Per due elementi, la probabilità che vengano estratti entrambi è $k(k-1)/(n(n-1))$.
\end{spiegazione}

\subsection*{Esercizio 5: Somma di Bernoulli con numero di termini di Poisson}
Sia $N \sim \text{Poisson}(\lambda)$ indipendente da $X_i \sim B(p)$. Qual è la legge della somma parziale aleatoria $S_N = \sum_{i=1}^N X_i$? E la legge di $N - S_N$?
\begin{spiegazione}
Condizionando rispetto a $N$, otteniamo $S_N|N \sim \text{Bin}(N, p)$. Integrando, la mgf o le probabilità rivelano che \textbf{$S_N \sim \text{Poisson}(\lambda p)$} e \textbf{$N - S_N \sim \text{Poisson}(\lambda(1-p))$}.
Inoltre, sorprendentemente, $S_N$ e $N-S_N$ (i successi e i fallimenti totali) sono variabili \textbf{indipendenti}! Questo si verifica calcolando la mgf congiunta.
\end{spiegazione}

\subsection*{Esercizio 6: Strategia del raddoppio al gioco della moneta}
Se perdo raddoppio la puntata, se vinco mi fermo. Si inizia puntando 1. Se il capitale iniziale è 1023, mi rovino se perdo 10 volte di fila ($2^{10}-1 = 1023$).
Il guadagno netto, se si vince entro i 10 turni, è sempre esattamente $+1$. Se si perde, è $-1023$.
\begin{spiegazione}
La durata del gioco $T$ è ristretta a $\{1, \dots, 10\}$. Si ha $\PP(T=n) = q^{n-1}p$ per $n \le 9$, e $\PP(T=10) = q^9$.
Il valore atteso di $T$ si calcola con somme finite di serie geometriche derivate.
Questo esercizio ci ricorda che le "strategie infallibili" come il raddoppio delle puntate collassano di fronte alla finitezza dei capitali, portando a rischi di rovina asimmetrici enormi.
\end{spiegazione}



\chapter{Processi stocastici a tempo discreto}

Introduciamo il concetto di processo stocastico a tempo discreto. I processi stocastici sono oggetti matematici il cui scopo è descrivere l'evoluzione temporale di fenomeni aleatori. 
Senza perdita di generalità, l'indice temporale sarà identificato con un sottoinsieme dei numeri naturali, solitamente a partire da $n=0$. Indicheremo l'insieme dei tempi con $\N$ (orizzonte finito) oppure $\N_0$ (orizzonte infinito). Nelle formule useremo $\N_0 = \N \cup \{0\}$.

\section{Nozione di Processo Stocastico}

Assumiamo fissato uno spazio di probabilità $(\Omega, \F, \PP)$ e uno spazio di arrivo $E$ (al più numerabile), munito della $\sigma$-algebra discreta (l'insieme delle parti).

\begin{definition}[Processo Stocastico]
Un processo stocastico a valori in $E$ è una famiglia $X = (X_n)_{n \in \N_0}$ di variabili aleatorie $X_n : \Omega \to E$.
\end{definition}

Possiamo pensare a $X$ come alla realizzazione in sequenza di variabili aleatorie nel tempo. Se fissiamo un esito $\omega \in \Omega$, la sequenza numerica (o a valori in $E$):
\[ (X_0(\omega), X_1(\omega), X_2(\omega), \dots) \]
prende il nome di \emph{traiettoria} del processo stocastico $X$.

\section{Filtrazioni e Struttura dell'Informazione}

Vogliamo dotare il modello di una struttura di informazione che rifletta l'apprendimento progressivo nel tempo: man mano che il tempo scorre, l'individuo "legge" cosa succede.
Supponiamo che l'individuo osservi i valori di un processo $\xi = (\xi_n)$. Al tempo $n$, egli conosce i valori $\xi_0(\omega), \dots, \xi_n(\omega)$. Questa conoscenza pregressa può modificare le aspettative sul futuro.
Formalizziamo questa idea con le \emph{filtrazioni}.

\begin{definition}[Filtrazione Naturale]
Sia $\xi$ un processo stocastico. Definiamo:
\[ \F_n^\xi = \sigma(\xi_0, \xi_1, \dots, \xi_n) \]
cioè la più piccola $\sigma$-algebra rispetto a cui tutte le variabili $\xi_0, \dots, \xi_n$ sono misurabili.
Chiaramente $\F_n^\xi \subseteq \F_{n+1}^\xi$. La famiglia crescente $(\F_n^\xi)_{n \in \N_0}$ si chiama \emph{filtrazione generata} da $\xi$ o \emph{filtrazione naturale}.
\end{definition}

\begin{definition}[Processo Adattato]
Un processo stocastico $X = (X_n)$ si dice \emph{adattato} alla filtrazione $(\F_n)$ se, per ogni tempo $n$, la variabile aleatoria $X_n$ è misurabile rispetto a $\F_n$.
\end{definition}

In parole povere, $X$ è adattato se il suo valore al tempo $n$ dipende unicamente dall'informazione disponibile fino al tempo $n$, senza sbirciare nel futuro. 
Spesso si usa la filtrazione naturale generata da $X$ stesso: in tal caso, il processo è banalmente adattato.

\section{Tempi di Arresto (Stopping Times)}

Una nozione cruciale per lo studio dei processi stocastici è quella di \emph{tempo di arresto} (stopping time). Formalizza l'idea di una regola d'arresto che si basa unicamente sull'informazione presente e passata, mai su quella futura.

\begin{definition}[Tempo di arresto]
Sia $(\F_n)_{n \in \N_0}$ una filtrazione. Una variabile aleatoria $\tau : \Omega \to \N_0 \cup \{\infty\}$ è detta \emph{tempo di arresto} se, per ogni $n \in \N_0$:
\[ \{\tau \le n\} \in \F_n \]
\end{definition}

\begin{hint}[Condizione equivalente]
Dire che $\{\tau \le n\} \in \F_n$ equivale a dire $\{\tau = n\} \in \F_n$.
A parole: l'informazione disponibile al tempo $n$ è sufficiente per stabilire inequivocabilmente se l'evento in questione "è già accaduto" o no.
\end{hint}

Un classico esempio di tempo di arresto è il \emph{tempo di primo passaggio} (hitting time).

\begin{definition}[Tempo di primo passaggio]
Sia $X$ un processo adattato e sia $A \subseteq E$. Il tempo di primo ingresso (o hitting time) in $A$ è definito come:
\[ \tau_A = \inf\{n \in \N_0 : X_n \in A\} \]
(con la convenzione $\inf \emptyset = \infty$).
\end{definition}

\begin{esempio}[Cosa \textbf{non} è un tempo di arresto]
Il tempo in cui un processo su orizzonte limitato $N$ raggiunge il suo \emph{massimo globale}, cioè $\tau = \operatorname{argmax}_{k \le N} X_k$, \textbf{non} è in generale un tempo di arresto. Per stabilire se $\tau = n$, bisognerebbe sapere se nei tempi futuri $n+1, \dots, N$ ci sarà un valore ancora più alto, ma questo richiede di "sbirciare nel futuro" (non è in $\F_n$).
\end{esempio}

\begin{definition}[Processo arrestato]
Dato un processo adattato $X_n$ e un tempo di arresto $\tau$, il \emph{processo arrestato} (stopped process) a $\tau$ è definito da:
\[ X_n^\tau = X_{\min(n, \tau)} = \begin{cases}
X_n & \text{se } n < \tau \\
X_\tau & \text{se } n \ge \tau
\end{cases} \]
\end{definition}

\begin{theorem}[Il processo arrestato è adattato]
Se $X$ è adattato alla filtrazione $(\F_n)$ e $\tau$ è un tempo di arresto, allora anche il processo arrestato $X^\tau$ è adattato alla medesima filtrazione.
\end{theorem}

\section{Esercizi Selezionati (Capitolo 3)}

\subsection*{Esercizio 1: Minimo di tempi di arresto}
Siano $\tau_1$ e $\tau_2$ due tempi di arresto. Dimostrare che $\tau = \min(\tau_1, \tau_2)$ è anch'esso un tempo di arresto.
\begin{spiegazione}
Per definizione, vogliamo verificare se $\{\tau \le n\} \in \F_n$. Osserviamo che il minimo tra due numeri è $\le n$ se e solo se almeno uno dei due è $\le n$:
\[ \{\min(\tau_1, \tau_2) \le n\} = \{\tau_1 \le n\} \cup \{\tau_2 \le n\} \]
Poiché entrambi gli eventi appartengono a $\F_n$ e $\F_n$ è una $\sigma$-algebra (chiusa rispetto alle unioni), l'unione appartiene a $\F_n$. Quindi $\tau$ è un tempo di arresto.
\end{spiegazione}

\subsection*{Esercizio 2: Tempo di ultima uscita (Last exit time)}
Sia $L_A = \sup\{n \in \N_0 : X_n \in A\}$ il tempo dell'ultima visita di un processo all'insieme $A$. Perché in generale \emph{non} è un tempo di arresto?
\begin{spiegazione}
Per determinare se $\{L_A = n\}$ si è verificato, non basta sapere che $X_n \in A$, ma occorre anche avere la certezza che per ogni $k > n$ il processo non tornerà in $A$ ($X_k \notin A$). Ma questa è un'informazione sul futuro, non contenuta in $\F_n$. Dunque non è una regola d'arresto (non puoi "fermarti" a $n$ sapendo che è l'ultima volta che tocchi $A$, perché non lo saprai finché non finirà il tempo!).
\end{spiegazione}

\subsection*{Esercizio 3: Atomi di una filtrazione}
Se $\xi_n \sim B(p)$ (lancio di moneta i.i.d.), descrivere gli atomi della $\sigma$-algebra $\F_1 = \sigma(\xi_0, \xi_1)$.
\begin{spiegazione}
Gli atomi (eventi elementari minimali) sono tutte le possibili realizzazioni congiunte delle prime due variabili, che descrivono completamente lo stato di conoscenza a tempo $n=1$. Poiché ciascuna può assumere 2 valori, vi sono $2^2 = 4$ atomi:
$A_1 = \{\xi_0=0, \xi_1=0\}$, $A_2 = \{\xi_0=0, \xi_1=1\}$, $A_3 = \{\xi_0=1, \xi_1=0\}$, $A_4 = \{\xi_0=1, \xi_1=1\}$.
Qualsiasi evento misurabile $\F_1$ è unione di questi atomi.
\end{spiegazione}

\subsection*{Esercizio 4: Somma di tempi di arresto}
Siano $\tau_1, \tau_2$ tempi di arresto. Mostrare che $\tau_1 + \tau_2$ è un tempo di arresto.
\begin{spiegazione}
L'evento $\{\tau_1 + \tau_2 = n\}$ si può scomporre come unione di tutti i modi in cui i due tempi possono sommare a $n$:
\[ \{\tau_1 + \tau_2 = n\} = \bigcup_{k=0}^n \big(\{\tau_1 = k\} \cap \{\tau_2 = n-k\}\big) \]
Per $k \le n$, l'evento $\{\tau_1 = k\} \in \F_k \subseteq \F_n$. Anche l'evento per $\tau_2$ appartiene a $\F_{n-k} \subseteq \F_n$. L'intersezione e la loro unione finita sono in $\F_n$. Dunque la somma è un tempo di arresto.
\end{spiegazione}



\chapter{Martingale}

In questo capitolo assumiamo fissato uno spazio di probabilità di riferimento $(\Omega, \F, \PP)$. A meno che non sia specificato altrimenti, $\xi = (\xi_n)_{n \in \N}$ sarà un processo stocastico e $(\F_n)_{n \in \N_0}$ la filtrazione generata da $\xi$.

\section{Nozione di Martingala}

Introduciamo una nozione importantissima nella teoria dei processi stocastici, utilizzata per descrivere giochi "equi".

\begin{definition}[Martingala, Submartingala, Supermartingala]
Sia $X = (X_n)_{n \in \N_0}$ un processo stocastico a valori reali, adattato a $(\F_n)$. Si dice che $X$ è una \emph{submartingala} (rispetto alla filtrazione data e alla probabilità $\PP$) se:
\begin{enumerate}
    \item $X_n \in L^1$ per ogni $n$ (ha valore atteso finito);
    \item $\E[X_{n+1} | \F_n] \ge X_n$ quasi certamente, per ogni $n$.
\end{enumerate}
Si dice \emph{supermartingala} se la disuguaglianza è inversa ($\le$). \\
Si dice \emph{martingala} se vale l'uguaglianza:
\[ \E[X_{n+1} | \F_n] = X_n \quad \text{q.c.} \]
\end{definition}

\begin{spiegazione}
Immagina che $X_n$ sia il tuo capitale al tempo $n$ in un gioco d'azzardo. 
Se $X$ è una \emph{martingala}, condizionatamente a tutto quello che sai fino ad oggi ($\F_n$), ti aspetti che domani il tuo capitale sia esattamente quello di oggi: stai facendo un \textbf{gioco equo}.
Se è una \emph{submartingala}, ti aspetti di avere di più: il gioco è \textbf{favorevole} (il prefisso sub- si riferisce al fatto che sta \emph{sotto} al valore futuro atteso).
Se è una \emph{supermartingala}, ti aspetti di avere di meno: il gioco è \textbf{sfavorevole} (come al casinò).
\end{spiegazione}

\subsection*{Esempi fondamentali}

\begin{esempio}[Passeggiata aleatoria]
Sia $\xi_n \in L^1$ una successione di variabili i.i.d. con $\E[\xi_n] = 0$. Definiamo $X_0 \in L^1$ e $X_n = X_{n-1} + \xi_n$. Allora:
\[ \E[X_n | \F_{n-1}] = \E[X_{n-1} + \xi_n | \F_{n-1}] = X_{n-1} + \E[\xi_n] = X_{n-1} \]
Questo processo (la somma di incrementi a media nulla) è una martingala.
\end{esempio}

\begin{esempio}[Martingala geometrica]
Sia $\xi_n$ una successione i.i.d. con $\E[\xi_n] = 1$. Se definiamo il processo moltiplicativo $X_n = X_{n-1} \cdot \xi_n$, allora:
\[ \E[X_n | \F_{n-1}] = X_{n-1} \E[\xi_n] = X_{n-1} \]
Anche questa è una martingala (modello base per la finanza).
\end{esempio}

\begin{esempio}[Martingala chiusa]
Sia $Z \in L^1$ una variabile aleatoria (il risultato di un esperimento che avverrà alla fine del tempo). Se definiamo la nostra stima ad ogni istante come $X_n = \E[Z | \F_n]$, allora per la Tower Property:
\[ \E[X_{n+1} | \F_n] = \E[\E[Z | \F_{n+1}] | \F_n] = \E[Z | \F_n] = X_n \]
$X_n$ è una martingala. Rappresenta l'aggiornamento dell'informazione su un target fisso.
\end{esempio}

\begin{theorem}[Media di una martingala]
Se $X$ è una submartingala, la funzione $n \mapsto \E[X_n]$ è non decrescente.
Se è supermartingala, è non crescente.
Se è una martingala, $\E[X_n]$ è \textbf{costante} nel tempo.
\end{theorem}
\begin{hint}[Condizione non sufficiente]
Avere la media costante \emph{non} basta per essere una martingala! La proprietà di martingala è un'uguaglianza condizionata (vale percorso per percorso, in media rispetto agli scenari futuri partendo da un punto preciso), non solo un'uguaglianza numerica assoluta. Un processo può avere media globale costante ma oscillare in modo prevedibile.
\end{hint}

\section{Risultati Utili sulle Martingale}

Alcuni risultati algebrici e teorici sono strumenti di lavoro quotidiano nella teoria.

\begin{theorem}[Compensatore Prevedibile]
Sia $X$ una martingala a quadrato integrabile ($\E[X_n^2] < \infty$). Allora:
\begin{enumerate}
    \item $X^2$ è una \emph{submartingala}. (Applicazione diretta della Disuguaglianza di Jensen condizionata per $\phi(x)=x^2$).
    \item Esiste un processo $C_n = \sum_{k=1}^n \E[(X_k - X_{k-1})^2 | \F_{k-1}]$, chiamato \emph{compensatore prevedibile}, tale che:
    \[ M_n = X_n^2 - C_n \]
    è una \textbf{martingala}. (Nota: prevedibile significa che $C_n$ è misurabile già al tempo $n-1$).
\end{enumerate}
\end{theorem}

\begin{theorem}[Trasformata di Martingala (Strategie di trading)]
Sia $X$ una martingala (i prezzi di mercato), e sia $h = (h_n)_{n \in \N}$ un processo limitato e \emph{prevedibile} (cioè $h_n$ è $\F_{n-1}$-misurabile, ovvero decido quanto investire $h_n$ basandomi sull'informazione di ieri).
Il guadagno cumulato $Y$ (o trasformata di martingala) è:
\[ Y_n = Y_0 + \sum_{k=1}^n h_k (X_k - X_{k-1}) \]
Allora \textbf{$Y$ è una martingala}.
\end{theorem}
\begin{spiegazione}
Significato profondo: non esiste alcuna strategia di scommessa (anche astuta e basata sul passato come il raddoppio), finché non sbircia il futuro, che possa trasformare un gioco equo (martingala $X$) in un gioco favorevole (submartingala $Y$). Il capitale risulterà sempre una martingala (gioco equo cumulato).
\end{spiegazione}

\begin{theorem}[Martingala arrestata (Stopped Martingale)]
Sia $X$ una martingala e $\tau$ un tempo di arresto. Il processo arrestato $X_n^\tau = X_{\min(n, \tau)}$ è anch'esso una martingala.
\end{theorem}
\begin{hint}[Riflessione strategica]
Questo teorema è un corollario della trasformata di martingala ponendo $h_n = \mathbf{1}_{\{\tau \ge n\}}$.
Il processo $h_n$ vale 1 finché non decidi di fermarti, poi vale 0. È ovviamente prevedibile (se ti sei fermato prima di oggi, lo sai già ieri). Di conseguenza, "smettere di giocare a un tempo di arresto" non altera l'equità del gioco.
\end{hint}



\chapter{Passeggiate Aleatorie (Random Walks)}

In questo capitolo esploriamo il comportamento a lungo termine di una particolare classe di processi: le \emph{passeggiate aleatorie}. 
Sia $\xi = (\xi_n)_{n \in \N}$ una successione di variabili aleatorie discrete a valori reali, e definiamo:
\[ X_0 = \xi_0, \quad X_n = X_{n-1} + \xi_n = \sum_{k=0}^n \xi_k \]
$X$ è un processo stocastico che rappresenta la posizione di una particella dopo $n$ passi di ampiezza $\xi_k$.

\section{Legge 0-1 di Kolmogorov e Lemmi di Borel-Cantelli}

Per studiare il comportamento asintotico ($n \to \infty$), è utile introdurre l'informazione portata dalla "coda" della successione.

\begin{definition}[$\sigma$-algebra coda]
Sia $\xi$ una successione di variabili. La \emph{$\sigma$-algebra coda} (tail $\sigma$-algebra) $\mathcal{T}^\xi$ è definita come l'intersezione di tutte le $\sigma$-algebre future:
\[ \mathcal{T}^\xi = \bigcap_{n \in \N} \sigma(\xi_n, \xi_{n+1}, \dots) \]
\end{definition}

Eventi che non dipendono da un numero finito di esiti iniziali (es. "la successione converge", "la somma diverge", "visita lo zero infinite volte") appartengono alla $\sigma$-algebra coda.

\begin{theorem}[Legge 0-1 di Kolmogorov]
Se le variabili $\xi_n$ sono \textbf{indipendenti}, allora ogni evento $A \in \mathcal{T}^\xi$ ha probabilità estrema:
\[ \PP(A) = 0 \quad \text{oppure} \quad \PP(A) = 1 \]
Di conseguenza, ogni variabile aleatoria misurabile rispetto a $\mathcal{T}^\xi$ è quasi certamente costante.
\end{theorem}

Per stabilire quale sia questa probabilità (0 o 1) per eventi del tipo "un evento si verifica infinite volte" (limsup di eventi), usiamo i lemmi di Borel-Cantelli.
Siano $(A_n)_{n \in \N}$ eventi e definiamo l'evento $\{A_n \text{ i.o.}\} = \limsup_{n \to \infty} A_n = \bigcap_{n} \bigcup_{k \ge n} A_k$ (i.o. sta per "infinitely often").

\begin{theorem}[Lemmi di Borel-Cantelli]
\textbf{Primo Lemma (Generale):} Se la serie delle probabilità converge:
\[ \sum_{n=1}^\infty \PP(A_n) < \infty \implies \PP(A_n \text{ i.o.}) = 0 \]
(Se un evento è "raro" in media, accadrà solo un numero finito di volte).

\textbf{Secondo Lemma (Richiede Indipendenza):} Se gli eventi $A_n$ sono indipendenti e la serie diverge:
\[ \sum_{n=1}^\infty \PP(A_n) = \infty \implies \PP(A_n \text{ i.o.}) = 1 \]
(Se tenti all'infinito un evento, anche con probabilità molto bassa, ma tale che la serie diverge, prima o poi capiterà infinite volte).
\end{theorem}

\begin{esempio}[Convergenza in probabilità ma non q.c.]
Sia $X_n$ una successione di Bernoulli indipendenti con $\PP(X_n = 1) = 1/n$.
Fissato $\epsilon \in (0, 1)$, $\PP(|X_n| > \epsilon) = 1/n \to 0$. Quindi $X_n \stackrel{P}{\to} 0$ in probabilità.
Tuttavia, $\sum \PP(X_n = 1) = \sum 1/n = \infty$. Essendo indipendenti, per il Secondo Lemma di Borel-Cantelli, $X_n = 1$ infinite volte quasi certamente. Quindi $X_n$ \emph{non converge} a 0 quasi certamente.
\end{esempio}

\section{Comportamento asintotico delle passeggiate simmetriche}

Concentriamoci su un caso specifico (modello base): i passi $\xi_n \in \{-1, 0, 1\}$ sono indipendenti e simmetrici, con $\PP(\xi_n = 1) = \PP(\xi_n = -1) = p_n \le 1/2$.

Cosa fa la particella nel lungo termine?
Usando Borel-Cantelli otteniamo la seguente dicotomia:
\begin{enumerate}
    \item Se $\sum_{n=1}^\infty p_n < \infty$, il primo lemma di Borel-Cantelli ci dice che la particella prima o poi "si ferma": da un certo $n$ in poi tutti i passi saranno nulli. (La passeggiata è definitivamente costante q.c.).
    \item Se $\sum_{n=1}^\infty p_n = \infty$, la passeggiata non si fermerà mai definitivamente. In realtà, si può dimostrare usando le martingale che in questo caso le fluttuazioni sono \textbf{illimitate}:
    \[ \limsup_{n \to \infty} X_n = +\infty \quad \text{e} \quad \liminf_{n \to \infty} X_n = -\infty \quad \text{q.c.} \]
\end{enumerate}

\section{Intermezzo: La rovina del giocatore (caso equo)}

Un'applicazione classica è la Rovina del Giocatore. 
Due giocatori (A) e (B) con capitali inziali $a, b \in \N$ scommettono 1 moneta a testa con esiti indipendenti (vince A, vince B o pareggio). Il gioco è \emph{equo}: la probabilità che vinca A in un round è uguale a quella che vinca B (chiamiamola $p > 0$).
Il gioco termina quando uno dei due finisce i soldi (capitale 0). Qual è la probabilità di rovina per A o per B?

Costruiamo una passeggiata aleatoria $X_n$ per il guadagno netto di (A). $X$ è una \emph{martingala} (poiché $\E[\xi] = p - p = 0$). Il gioco finisce al tempo d'arresto $\tau = \inf\{n : X_n = -a \text{ oppure } X_n = b\}$.
Sappiamo che la passeggiata in un gioco infinito ($\sum p = \infty$) oscilla illimitatamente, quindi $\tau < \infty$ quasi certamente (il gioco finirà sicuramente).

Applicando il teorema della martingala arrestata:
\[ \E[X_{\tau \wedge n}] = \E[X_0] = 0 \]
Usando la convergenza dominata (poiché il processo arrestato è confinato tra $-a$ e $b$), si ottiene $\E[X_\tau] = 0$.
Ma $X_\tau$ può valere solo $-a$ (se (A) si rovina, con prob. $\PP(R_A)$) oppure $b$ (se (B) si rovina, con prob. $\PP(R_B)$).
Quindi:
\[ -a \PP(R_A) + b \PP(R_B) = 0 \quad \text{e sapendo che} \quad \PP(R_A) + \PP(R_B) = 1 \]
Otteniamo la formula classica:
\[ \PP(R_A) = \frac{b}{a+b}, \quad \PP(R_B) = \frac{a}{a+b} \]

\begin{hint}[Morale]
In un gioco equo, la probabilità di sopravvivenza è strettamente proporzionale al capitale iniziale. Il giocatore con le "tasche più fonde" ha enormemente più probabilità di mandare in rovina l'avversario.
Inoltre, siccome in questo caso la passeggiata raggiunge q.c. un qualsiasi livello $\pm \alpha$, se (B) avesse un "capitale infinito" (come un Casinò), (A) andrebbe inevitabilmente in rovina con probabilità 1: $\lim_{b \to \infty} \PP(R_A) = 1$.
\end{hint}



\chapter{Modelli Finanziari in Spazi di Probabilità Finiti}

In questo capitolo esploriamo i fondamenti della matematica finanziaria in tempo discreto. L'obiettivo principale è mostrare come la teoria delle martingale costituisca lo strumento fondamentale per prezzare (valutare) i derivati finanziari e descrivere l'assenza di arbitraggio.

\section{Modelli uni-periodali (One-period models)}

Consideriamo un mercato finanziario elementare che vive solo in due istanti di tempo: $n=0$ (oggi) e $n=1$ (domani). Lo spazio di probabilità $(\Omega, \mathcal{F}, \PP)$ è finito: $\Omega = \{\omega_1, \dots, \omega_M\}$, con probabilità $\PP(\{\omega_k\}) = p_k > 0$.

Il mercato è composto da due asset:
\begin{itemize}
    \item Un asset privo di rischio (\textbf{Bond} o Conto in banca) $B$:
    \[ B_0 = 1, \quad B_1 = 1 + r \]
    dove $r > -1$ è il tasso di interesse (deterministico).
    \item Un asset rischioso (\textbf{Stock} o Azione) $S$:
    \[ S_0 = s > 0, \quad S_1(\omega) = s_k > 0 \text{ se } \omega = \omega_k \]
    Il valore $S_1$ è una variabile aleatoria (il prezzo di domani non è noto oggi).
\end{itemize}

È comodo introdurre i \emph{valori scontati} degli asset (prendendo il bond come \textit{numeraire}):
\[ B^*_n = \frac{B_n}{(1+r)^n} \implies B^*_0 = 1, \; B^*_1 = 1 \]
\[ S^*_n = \frac{S_n}{(1+r)^n} \implies S^*_0 = s, \; S^*_1 = \frac{S_1}{1+r} \]

\subsection{Portafogli e Arbitraggio}

Una \textbf{strategia di portafoglio} è un vettore $h = (h^B, h^S) \in \R^2$ che indica le quantità di bond e stock possedute. Il valore del portafoglio agli istanti $n=0, 1$ è:
\[ V_0^h = h^B B_0 + h^S S_0 = h^B + h^S s \]
\[ V_1^h = h^B B_1 + h^S S_1 = h^B (1+r) + h^S S_1 \]
Anche per il portafoglio definiamo il valore scontato: $V_n^{*h} = V_n^h / (1+r)^n = h^B B^*_n + h^S S^*_n$.

\begin{definition}[Arbitraggio]
Un \textbf{arbitraggio} è una strategia di portafoglio $h$ tale che:
\begin{enumerate}
    \item $V_0^h = 0$ (investimento iniziale nullo);
    \item $\PP(V_1^h \ge 0) = 1$ (nessun rischio di perdere denaro);
    \item $\PP(V_1^h > 0) > 0$ (probabilità strettamente positiva di fare profitto).
\end{enumerate}
Se non esistono strategie con queste proprietà, il mercato si dice \textbf{non arbitraggiabile} (arbitrage-free).
\end{definition}

\begin{hint}[Spiegazione Semplice dell'Arbitraggio]
Un arbitraggio è, letteralmente, una "macchina per fare soldi gratis". Entri nel mercato con zero euro (ad esempio andando corto sul bond per finanziare l'acquisto di stock), sei certo di non perdere mai nulla e hai una reale possibilità di uscirne più ricco. In un mercato efficiente, simili opportunità vengono immediatamente sfruttate (e quindi distrutte) dagli investitori. Pertanto, l'assenza di arbitraggio è il requisito minimo per un modello finanziario sensato.
\end{hint}

\subsection{Misure Martingala e Teoremi Fondamentali}

Per collegare i prezzi all'assenza di arbitraggio, introduciamo il concetto chiave del capitolo.

\begin{definition}[Misura Martingala Equivalente]
Una probabilità $\QQ$ su $\Omega$ si dice \textbf{equivalente} a $\PP$ (scritto $\QQ \sim \PP$) se hanno gli stessi eventi impossibili ($\QQ(A)=0 \iff \PP(A)=0$).
Una misura equivalente $\QQ$ è detta \textbf{Misura Martingala Equivalente (EMM)} se il prezzo scontato dello stock $S^*$ è una martingala sotto $\QQ$:
\[ \E^\QQ[S^*_1] = S^*_0 \iff \E^\QQ[S_1] = s(1+r) \]
\end{definition}

Sotto la probabilità $\QQ$, il rendimento atteso dell'asset rischioso è pari al rendimento dell'asset privo di rischio. $\QQ$ è nota anche come probabilità \emph{risk-neutral}.

\begin{theorem}[Primo Teorema Fondamentale dell'Asset Pricing]
Il mercato è non arbitraggiabile (arbitrage-free) \textbf{se e solo se} esiste (almeno) una misura martingala equivalente $\QQ$.
\end{theorem}

Un'altra questione fondamentale è la "completezza" del mercato. Un \textbf{Contingent Claim} (es. un'opzione) è un contratto che paga $\phi(S_1)$ al tempo 1. Il claim è detto \textbf{replicabile (o hedgeable)} se esiste un portafoglio $h$ tale che $V_1^h = \phi(S_1)$ per ogni possibile stato $\omega$. 
Se \emph{tutti} i possibili claim sono replicabili, il mercato si dice \textbf{completo}.

\begin{theorem}[Secondo Teorema Fondamentale dell'Asset Pricing]
Assumendo che il mercato sia non arbitraggiabile, esso è \textbf{completo} se e solo se esiste \textbf{esattamente una} misura martingala equivalente $\QQ$.
\end{theorem}

In questo modello a singolo stock, se $S_1$ può assumere 2 valori ($M=2$, modello binomiale) il mercato è completo. Se assume $M > 2$ valori, il mercato è incompleto (ci sono infinite misure $\QQ$).

\section{Pricing (Valutazione dei derivati)}

Se consideriamo un claim $\phi(S_1)$ in un mercato arbitrage-free, qual è il suo "prezzo giusto" (o \textit{fair price}) $x$ al tempo $n=0$?
Il prezzo $x$ deve essere tale che se inseriamo il derivato nel mercato come nuovo asset da scambiare, il mercato allargato rimanga privo di arbitraggi.

Abbiamo due strade (che coincidono per claim replicabili):
\begin{enumerate}
    \item \textbf{Hedging (Replica):} Se il claim è replicabile dal portafoglio $h$, allora per evitare arbitraggi il prezzo \emph{deve} essere il costo di replica oggi:
    \[ x = V_0^h \]
    Se fosse $x > V_0^h$, venderemmo il claim e compreremmo il portafoglio $h$, incassando la differenza subito a rischio zero. Viceversa se $x < V_0^h$.
    \item \textbf{Risk-Neutral Pricing:} In base al Primo Teorema, per non avere arbitraggi, anche il prezzo scontato del derivato deve essere una martingala sotto $\QQ$:
    \[ x = \E^\QQ\left[ \frac{\phi(S_1)}{1+r} \right] \]
\end{enumerate}

\begin{tcolorbox}[colback=red!5!white,colframe=red!75!black,title=Attenzione: Pricing sotto $\PP$ vs $\QQ$]
L'errore più comune nei modelli finanziari è calcolare il prezzo di un'opzione scontando il suo valore atteso rispetto alla probabilità oggettiva (o fisica) $\PP$. \textbf{Questo è sbagliato e porta ad arbitraggi!} Il prezzo di non-arbitraggio \emph{deve} essere calcolato usando il valore atteso sotto la misura risk-neutral $\QQ$. Sotto $\QQ$ gli investitori non richiedono alcun premio per il rischio.
\end{tcolorbox}

\section{Modelli multi-periodali: Il modello Binomiale (CRR)}

Estendiamo l'orizzonte a $N$ periodi di tempo: $n=0, 1, \dots, N$. Il mercato ha un bond deterministico $B_n = (1+r)^n$. L'asset rischioso $S_n$ evolve stocasticamente.

Una strategia $h = (h_n^B, h_n^S)_{n=0}^{N}$ è ora un \emph{processo pre-vedibile} (decido $h_n$ all'istante $n-1$, in base alle informazioni disponibili). La strategia è \textbf{autofinanziante} (self-financing) se le variazioni del valore di portafoglio sono dovute unicamente all'andamento del mercato, senza immissioni/prelievi di liquidità:
\[ h_{n+1}^B B_n + h_{n+1}^S S_n = h_n^B B_n + h_n^S S_n = V_n^h \]
Passando ai valori scontati:
\[ V_{n+1}^{*h} - V_{n}^{*h} = h_{n+1}^S (S_{n+1}^* - S_n^*) \]
che è esattamente una trasformata di martingala discreta! Ecco perché se $S^*$ è martingala (sotto $\QQ$), anche l'evoluzione di qualsiasi portafoglio autofinanziato $V^{*h}$ sarà una martingala sotto $\QQ$.

\subsection{Il Modello Binomiale di Cox-Ross-Rubinstein}
Nel modello binomiale, il prezzo dello stock passa dall'istante $n$ a $n+1$ moltiplicandosi per un fattore $Z_{n+1}$ che può assumere solo due valori: $u$ (up) e $d$ (down).
\[ S_{n+1} = Z_{n+1} S_n, \quad Z_{n+1} \in \{u, d\}, \quad 0 < d < u \]
Per evitare arbitraggi banali (in cui $S$ rende sempre più o sempre meno del bond), imponiamo:
\[ d < 1+r < u \]

Sotto questa ipotesi, l'equazione per la martingalità del prezzo scontato, ovvero $\E^\QQ[Z_{n+1}] = 1+r$, ci forza a determinare univocamente le probabilità risk-neutral di salita e discesa in ogni singolo passo temporale:
\[ q_u = \QQ(Z = u) = \frac{(1+r) - d}{u - d} \in (0, 1) \]
\[ q_d = \QQ(Z = d) = \frac{u - (1+r)}{u - d} \in (0, 1) \]
Essendo la misura $\QQ$ unica, il Primo e Secondo Teorema garantiscono che il Modello Binomiale sia \textbf{arbitrage-free} e \textbf{completo}.

\subsection{Prezzatura backward}
Data un'opzione europea con payoff $\phi(S_N)$ a scadenza $N$, il suo prezzo ad un generico istante $n$ è dato dal valore atteso condizionato:
\[ X_n = \frac{1}{(1+r)^{N-n}} \E^\QQ[\phi(S_N) \mid \mathcal{F}_n] \]
Nel modello binomiale, questa formula si esprime con la \emph{backward recursion}. Al tempo finale $X_N = \phi(S_N)$. Per un passo all'indietro:
\[ X_n(s) = \frac{1}{1+r} \left( q_u X_{n+1}(su) + q_d X_{n+1}(sd) \right) \]
Costruendo l'albero recombining dei prezzi possibili (binomial tree), possiamo partire dalle foglie (istante $N$) e lavorare a ritroso (istante per istante) calcolando per ogni nodo il valore pesato $X_n$ fino ad arrivare a $X_0$, che sarà il fair price al tempo iniziale.



\chapter{Catene di Markov (Spazio Discreto)}

Le Catene di Markov sono una classe fondamentale di processi stocastici utilizzati per modellizzare fenomeni in cui l'evoluzione futura del sistema dipende unicamente dal suo stato presente, e non da come ci sia arrivato. Questa proprietà è nota come "assenza di memoria".

\section{Introduzione Euristica}

\begin{esempio}[Il meteo di Salisburgo]
A Salisburgo il meteo può essere Soleggiato (S), Piovoso (P) o Nevoso (N). Supponiamo valgano queste regole empiriche:
\begin{itemize}
    \item Non ci sono mai due giorni di sole di fila. Se oggi c'è sole, domani nevica o piove (50\% e 50\%).
    \item Se oggi piove o nevica, domani c'è una probabilità del 50\% che il tempo resti invariato, e 25\% per ciascuna delle altre due opzioni.
\end{itemize}
L'informazione su cosa farà domani è completamente determinata dal tempo di \emph{oggi}. Il fatto che ieri piovesse non aggiunge alcuna informazione utile. Questa è l'essenza della \textbf{proprietà di Markov}.
Le domande tipiche a cui risponde la teoria sono: "Qual è la probabilità che tra due settimane piova?" oppure "In media, quanti giorni di sole ci sono in un anno (comportamento a lungo termine)?"
\end{esempio}

\section{Definizione e Costruzione}

Sia $(\Omega, \F, \PP)$ uno spazio di probabilità e $(\F_n)_{n \in \N}$ la filtrazione di riferimento. Un processo stocastico $X = (X_n)_{n \in \N}$ a valori in uno spazio \emph{discreto} $E$ è una \textbf{Catena di Markov omogenea} se:
\[ \PP(X_{n+1} = y \mid \F_n) = \PP(X_{n+1} = y \mid X_n) = p(X_n, y) \quad \text{q.c.} \]
per ogni $y \in E$.

Gli ingredienti fondamentali per definire la dinamica del sistema sono due:
\begin{enumerate}
    \item La \textbf{distribuzione iniziale} $\mu_0$: la legge di $X_0$ (es. partiamo dallo stato $x$ con probabilità 1).
    \item Le \textbf{probabilità di transizione} $p(x, y)$: la probabilità di passare dallo stato $x$ allo stato $y$ in un passo. 
\end{enumerate}

Se lo spazio $E$ è finito ($|E| = N$), le probabilità di transizione possono essere racchiuse in una \textbf{Matrice di transizione} $M \in \R^{N \times N}$, in cui ogni riga somma a 1 (matrice stocastica):
\[ \sum_{y \in E} p(x, y) = 1 \quad \forall x \in E \]

\begin{theorem}[Esistenza e Unicità]
Dato uno spazio $E$, una distribuzione iniziale $\mu_0$ e un nucleo di transizione $p(x, y)$, esiste una Catena di Markov con tali caratteristiche, e la sua legge (cioè le probabilità su tutte le possibili traiettorie) è univocamente determinata.
\end{theorem}

\subsection{Dinamica a $n$ passi (Chapman-Kolmogorov)}

Se $p(x,y)$ descrive la mossa in 1 passo, indichiamo con $p^{(n)}(x,y)$ la probabilità di andare da $x$ a $y$ in $n$ passi. Sfruttando la proprietà di Markov, si ottiene l'equazione di \textbf{Chapman-Kolmogorov}:
\[ p^{(n+m)}(x, y) = \sum_{z \in E} p^{(m)}(x, z) p^{(n)}(z, y) \]
In forma matriciale (per $E$ finito), questo significa semplicemente che la matrice di transizione a $n$ passi è la potenza $n$-esima della matrice a 1 passo: $M^{(n)} = M^n$.

Un modo equivalente per guardare la dinamica è la \textbf{transizione di legge}. Se $\mu_n$ è il vettore colonna (o riga a seconda della convenzione) delle probabilità di trovarsi in ciascuno stato al tempo $n$, allora:
\[ \mu_{n+1} = \mu_n M \implies \mu_n = \mu_0 M^n \]

\section{La Proprietà di Markov Forte}

Un risultato molto profondo asserisce che la catena di Markov perde la memoria non solo agli istanti di tempo deterministici (es. $n=10$), ma anche agli istanti casuali descritti da \textbf{tempi d'arresto} $\tau$.

\begin{theorem}[Proprietà di Markov Forte]
Sia $\tau$ un tempo d'arresto. Condizionatamente all'evento $\{\tau < \infty\}$, il processo post-arresto:
\[ Y_k := X_{\tau + k} \]
è ancora una Catena di Markov con la \emph{stessa} matrice di transizione di $X$, indipendente dal passato osservato fino a $\tau$.
\end{theorem}

\section{Classificazione degli Stati}

\subsection{Comunicazione e Classi Irriducibili}
Diciamo che:
\begin{itemize}
    \item $y$ è \textbf{accessibile} da $x$ ($x \to y$) se c'è una probabilità positiva di arrivarci in un certo numero di passi ($p^{(n)}(x,y) > 0$).
    \item $x$ e $y$ \textbf{comunicano} ($x \leftrightarrow y$) se $x \to y$ e $y \to x$.
\end{itemize}
La comunicazione è una relazione di equivalenza. Partiziona $E$ in \textbf{classi irriducibili}. Se tutto $E$ forma un'unica classe, la catena si dice \textbf{Irriducibile} (si può viaggiare da qualsiasi stato a qualsiasi altro stato).

\subsection{Ricorrenza e Transitorietà}
Fissato lo stato di partenza $X_0 = x$, studiamo il tempo (eventuale) di primo ritorno in $x$:
\[ \tau_x := \inf\{n \ge 1 : X_n = x\} \]
Definiamo la probabilità di \emph{ritorno in $x$ in tempo finito}:
\[ F(x, x) = \PP_x(\tau_x < \infty) \]

\begin{definition}[Ricorrenza e Transitorietà]
Uno stato $x$ è:
\begin{itemize}
    \item \textbf{Ricorrente} se $F(x, x) = 1$ (è assolutamente certo che la particella tornerà in $x$).
    \item \textbf{Transiente} se $F(x, x) < 1$ (c'è una probabilità positiva che la particella se ne vada per sempre senza mai più tornare in $x$).
\end{itemize}
\end{definition}

Sia $V_x = \sum_{n=1}^\infty \1_{\{X_n = x\}}$ il numero totale di visite ad $x$. Si dimostra che:
\begin{itemize}
    \item Se $x$ è transiente, $V_x < \infty$ q.c. La particella \textbf{abbandona definitivamente} $x$. Criterio analitico: $\sum_n p^{(n)}(x,x) < \infty$.
    \item Se $x$ è ricorrente, $V_x = \infty$ q.c. La particella visiterà $x$ \textbf{infinite volte}. Criterio analitico: $\sum_n p^{(n)}(x,x) = \infty$.
\end{itemize}

\begin{hint}[Classificazione e Contagio]
In uno spazio finito non possono essere tutti stati transienti (sennò la particella "uscirebbe" dallo spazio, assurdo). 
Inoltre, la ricorrenza è "contagiosa": se $x$ è ricorrente e $x \to y$, allora anche $y$ deve essere ricorrente. Per questo, in una classe irriducibile, \textbf{gli stati sono o tutti ricorrenti o tutti transienti}. Se la catena è irriducibile ed $E$ è finito, essa è per forza Ricorrente.
\end{hint}

\section{Misure Invarianti e Teorema Ergodico}

\begin{definition}[Distribuzione Invariante]
Una distribuzione di probabilità $\mu^*$ su $E$ è \textbf{invariante} o stazionaria se per ogni $y$:
\[ \sum_{x \in E} \mu^*(x) p(x,y) = \mu^*(y) \iff \mu^* = \mu^* M \]
In termini di algebra, $\mu^*$ è un autovettore sinistro della matrice di transizione relativo all'autovalore 1. Se il sistema parte distribuito come $\mu^*$, la distribuzione rimarrà $\mu^*$ per tutti gli istanti successivi (equilibrio dinamico).
\end{definition}

Se uno stato è transiente, il suo "peso" all'equilibrio è zero: $\mu^*(x) = 0$.
Affinché esista una distribuzione invariante, la catena (se irriducibile) deve non solo tornare (ricorrente), ma \emph{tornare in fretta}. Questo introduce l'ultima classificazione per gli stati ricorrenti:

\begin{itemize}
    \item \textbf{Null-Ricorrente}: Ritorna al $100\%$ ($F=1$), ma il tempo medio di attesa per il ritorno è infinito ($\E_x[\tau_x] = \infty$). \emph{Non ammette misura di probabilità invariante}.
    \item \textbf{Positivo-Ricorrente}: Il tempo medio di ritorno è finito ($\E_x[\tau_x] < \infty$). \emph{Esiste una probabilità invariante, ed è proporzionale a:}
    \[ \mu^*(x) = \frac{1}{\E_x[\tau_x]} \]
\end{itemize}

\begin{spiegazione}[Misure stazionarie e tempi di ritorno]
C'è un profondo legame tra frequenza e probabilità limite. Se in media impiego 10 giorni a tornare a casa, significa che a lungo termine spenderò a casa 1 giorno su 10 (il 10\% del tempo). Ecco perché la probabilità stazionaria $\mu^*(x)$ è l'inverso del tempo medio di primo ritorno!
\end{spiegazione}

\begin{theorem}[Teorema Ergodico]
Sia $X$ una catena irriducibile e positivo-ricorrente (dunque con probabilità invariante $\mu^*$). Allora per ogni funzione $f$ ragionevole (integrabile rispetto a $\mu^*$), la media nel \textbf{tempo} converge alla media nello \textbf{spazio} quasi certamente:
\[ \lim_{n \to \infty} \frac{1}{n} \sum_{k=1}^n f(X_k) = \sum_{x \in E} f(x) \mu^*(x) \quad \text{q.c.} \]
\end{theorem}

Il Teorema Ergodico generalizza la Legge dei Grandi Numeri (LLN). La LLN richiede indipendenza; il teorema ergodico permette "memoria", a patto che la catena si mescoli bene (positivo-ricorrente) finendo per dimenticare le condizioni iniziali.

\section{Esercizi Tipici}

Un classico esercizio di esame (o progetto) richiede di, dato un piccolo grafo o una matrice $M$:
\begin{enumerate}
    \item Disegnare il grafo delle transizioni e identificare le \textbf{classi irriducibili}.
    \item Stabilire chi è \textbf{Transiente} (chi prima o poi finisce in una trappola e non torna) e chi è \textbf{Ricorrente}.
    \item Se lo spazio è finito, la classe ricorrente è automaticamente positivo-ricorrente. Si imposta quindi il sistema $\mu^* M = \mu^*$ unito a $\sum \mu^*(x) = 1$ per trovare l'unica \textbf{Distribuzione Invariante}.
    \item Trovare i tempi medi di ritorno calcolando $\E_x[\tau_x] = 1 / \mu^*(x)$.
    \item Usare il Teorema Ergodico per stimare integrali temporali a lungo termine.
\end{enumerate}



\chapter{Controllo Ottimo Stocastico a Tempo Discreto}

In questo capitolo introduciamo la teoria del controllo ottimo stocastico a tempo discreto e illustriamo il metodo della \textbf{Programmazione Dinamica} per risolverlo.

\section{Formulazione Generale del Problema}

Molti problemi applicati (finanza, economia, gestione scorte) rientrano in questo schema generale.
Gli elementi costitutivi sono:
\begin{enumerate}
    \item \textbf{Spazio degli stati} $S \subset \R^k$: dove vive il sistema.
    \item \textbf{Spazio dei controlli} $C \subset \R^d$: le azioni a disposizione dell'agente.
    \item Un \textbf{processo stocastico} $\xi = (\xi_n)_{n \in \N}$ i.i.d., che rappresenta l'alea del sistema (il caso).
    \item La \textbf{dinamica} del sistema, data da una funzione $f$:
    \[ X_{n+1} = f(X_n, c_n, \xi_{n+1}) \]
    dove $X_n$ è lo stato attuale, $c_n$ l'azione scelta e $\xi_{n+1}$ lo shock casuale in ingresso.
    \item I \textbf{vincoli} sulle azioni ammissibili $c_n \in \mathcal{A}(x)$.
    \item Un \textbf{payoff corrente} (funzione di running cost/reward) $g(X_n, c_n)$ e un \textbf{payoff finale} $\Phi(X_N)$.
\end{enumerate}

L'obiettivo è massimizzare (o minimizzare) il valore atteso della somma dei payoff, solitamente scontati con un fattore $\beta \in (0, 1]$:
\[ \sup_{c \in \mathcal{A}(x)} \E \left[ \sum_{n=0}^{N-1} \beta^n g(X_n, c_n) + \beta^N \Phi(X_N) \right] \]

\begin{spiegazione}[Anello Aperto vs Anello Chiuso]
Il controllo $c_n$ deve essere non-anticipante, ovvero deciso basandosi solo sulle informazioni disponibili fino al tempo $n$. 
Si dice "in anello aperto" se dipende dall'intera storia passata, o "in anello chiuso" (di feedback markoviano) se $c_n = \gamma(X_n)$, cioè l'azione ottimale dipende solo dallo stato \emph{corrente}.
\end{spiegazione}

\section{Esempi Classici}

\subsection{Crescita Economica Stocastica Ottima}
Un agente (o nazione) produce una ricchezza $Q_n$ ad ogni anno. Tale ricchezza dipende da un fattore di produttività casuale (es. clima, tecnologia) $Z_n$. Ad ogni anno, l'agente sceglie quanta parte della ricchezza $c_n$ consumare. La parte non consumata viene reinvestita e produce la ricchezza dell'anno successivo. L'obiettivo è massimizzare l'utilità attesa dei consumi nel tempo: $\E \sum \beta^n U(c_n)$.

\subsection{Allocazione di Portafoglio}
Hai un conto in banca a tasso privo di rischio $r$ e un asset rischioso (es. azioni) i cui ritorni aleatori sono $\mu_n$. Il tuo stato è il valore totale del portafoglio $W_n$. Il controllo $\pi_n$ è la frazione di ricchezza che investi nel titolo rischioso. L'obiettivo è massimizzare l'utilità della ricchezza finale: $\E [ U(W_N) ]$.

\subsection{Gestione Scorte (Inventory)}
Un magazzino deve rifornirsi ordinando $c_n$ unità di merce ad ogni intervallo. Esiste una domanda casuale $D_n$. C'è un costo di acquisto per $c_n$ e un costo di penalità se le scorte finali $Q_n$ si discostano troppo dalla domanda. L'obiettivo è minimizzare i costi attesi totali.

\section{Programmazione Dinamica (DP)}

L'idea geniale della Programmazione Dinamica (di Richard Bellman) è di "rompere" il problema su tutto l'orizzonte in una sequenza di problemi a singolo passo, procedendo a ritroso nel tempo.

\subsection{La Funzione Valore e l'Equazione di Bellman}

Definiamo la \textbf{Funzione Valore} $V_k(x)$ come il massimo payoff atteso che si può ottenere \emph{se al tempo $k$ ci troviamo nello stato $x$}.

\begin{theorem}[Principio di Ottimalità di Bellman]
La funzione valore soddisfa la seguente relazione ricorsiva (Equazione di Bellman):
\[ V_k(x) = \sup_{\gamma \in \Gamma(x)} \left\{ g(x, \gamma) + \beta \E \left[ V_{k+1}(f(x, \gamma, \xi_{k+1})) \right] \right\} \]
\end{theorem}

\begin{hint}[Come Leggere Bellman]
Il significato intuitivo è: il valore ottimale al tempo $k$ si ottiene scegliendo la mossa $\gamma$ che massimizza la somma tra il guadagno immediato oggi ($g(x, \gamma)$) e il valore atteso che avrò da domani in poi, considerando che mi ritroverò nel nuovo stato generato da $\gamma$ e dal caso.
\end{hint}

\subsection{Algoritmo in Orizzonte Finito}
Se il tempo $N$ è finito, il problema si risolve partendo dalla fine.
\begin{enumerate}
    \item Si pone $V_N(x) = \Phi(x)$ (al momento terminale, incasso semplicemente il payoff finale).
    \item Per $k = N-1, N-2, \dots, 0$, si calcola $V_k(x)$ applicando l'Equazione di Bellman, supponendo noto $V_{k+1}$.
    \item Al contempo, la mossa ottimale $\gamma^*_k(x)$ è quella che realizza il sup nell'equazione di Bellman.
\end{enumerate}

\subsection{Orizzonte Infinito e Teorema di Verifica}
Se $N = \infty$, il tempo scompare ($V_k(x) = V(x)$ per ogni $k$, dato che la dinamica è stazionaria). L'Equazione di Bellman diventa una equazione a punto fisso per una singola funzione $V(x)$:
\[ V(x) = \sup_{\gamma \in \Gamma(x)} \left\{ g(x, \gamma) + \beta \E \left[ V(f(x, \gamma, \xi_1)) \right] \right\} \]

Poiché procedere a ritroso dall'infinito è impossibile, in questo caso si usa un approccio \emph{Guess \& Verify}:
\begin{enumerate}
    \item "Indovino" una funzione $v_0(x)$ che risolve l'equazione di Bellman.
    \item Trovo la strategia di feedback ottima $\gamma^*(x)$.
    \item Verifico una condizione di \emph{trasversalità} (o crescita controllata all'infinito): $\lim_{N \to \infty} \E[v_0(X_N)] \le 0$ (o equazione simile, a seconda dei segni del payoff).
    \item Il \textbf{Teorema di Verifica} assicura allora che $v_0(x)$ è \emph{la} vera funzione valore e $\gamma^*$ è il vero controllo ottimo.
\end{enumerate}

\begin{esempio}[Esempio di Orizzonte Infinito]
Nel problema di crescita (Sect 8.4.1), se supponiamo che la funzione valore abbia la forma $v_0(x) = v_0(q, z) = \alpha(z) \ln q$, sostituendola nell'equazione di Bellman otteniamo i coefficienti ottimi, trovando così analiticamente sia il valore atteso massimo sia la porzione di ricchezza esatta da consumare ogni anno in base allo stato.
\end{esempio}



\appendix
\chapter{Richiami di Analisi e Serie}

Questa appendice contiene nozioni fondamentali di analisi matematica necessarie per lo studio dei modelli probabilistici, in particolare limiti superiore/inferiore e serie.

\section{Liminf e Limsup di successioni reali}

Sia $(x_n)_{n \in \N}$ una successione di numeri reali. 
La successione dell'estremo inferiore definitivo è definita come:
\[ y_n := \inf_{k \ge n} x_k \]
Questa successione $(y_n)$ è sempre crescente. Di conseguenza, ammette sempre limite (finito o infinito).

\begin{definition}[Limite Inferiore]
Il \textbf{limite inferiore} di una successione $(x_n)$ è definito come:
\[ \liminf_{n \to \infty} x_n := \lim_{n \to \infty} \left( \inf_{k \ge n} x_k \right) \]
\end{definition}

In maniera simmetrica, definendo la successione degli estremi superiori definitivi $z_n := \sup_{k \ge n} x_k$, che è sempre decrescente:

\begin{definition}[Limite Superiore]
Il \textbf{limite superiore} di una successione $(x_n)$ è definito come:
\[ \limsup_{n \to \infty} x_n := \lim_{n \to \infty} \left( \sup_{k \ge n} x_k \right) \]
\end{definition}

\begin{spiegazione}[A cosa servono $\liminf$ e $\limsup$?]
Il limite normale $\lim x_n$ potrebbe non esistere (pensa a $x_n = (-1)^n$, che oscilla tra -1 e 1 in eterno). Tuttavia, $\liminf$ e $\limsup$ esistono \textbf{sempre}, perché misurano i confini della successione a lungo termine. 
Se il limite normale esiste, allora $\liminf$ e $\limsup$ sono uguali, e coincidono col limite!
\end{spiegazione}

In generale, vale la disuguaglianza:
\[ \liminf_{n \to \infty} x_n \le \limsup_{n \to \infty} x_n \]
E si ha l'uguaglianza $\liminf x_n = \limsup x_n$ se e solo se il limite $\lim x_n$ esiste.

\section{Serie}

\subsection{Serie a termini non negativi}

Sia $I$ un insieme di indici (es. numerabile, con la stessa cardinalità di $\N$) e $f : I \to [0, +\infty]$. Data una biiezione $\sigma : \N \to I$, definiamo la somma parziale:
\[ S_n := \sum_{k=0}^n f(\sigma(k)) \]
Poiché $f \ge 0$, la successione $(S_n)$ è crescente, e perciò ammette limite (finito o infinito). Questo limite:
\[ \lim_{n \to \infty} \sum_{k=0}^n f(\sigma(k)) = \sum_{x \in I} f(x) \]
\textbf{non dipende} dalla scelta della biiezione $\sigma$. Ovvero, in una serie a termini positivi, puoi riordinare gli addendi come preferisci, il risultato finale sarà sempre lo stesso.

\subsection{Serie assolutamente convergenti}

Cosa succede se i termini possono essere negativi? Sia $f : I \to \R$.
Supponiamo che la serie dei valori assoluti converga:
\[ \sum_{x \in I} |f(x)| < \infty \]
In tal caso, la serie si dice \textbf{assolutamente convergente}. Anche in questo caso, è possibile dimostrare che il limite della somma:
\[ \sum_{x \in I} f(x) \]
esiste finito, e soprattutto \textbf{non dipende dall'ordine in cui si sommano i termini}.

\begin{hint}[Attenzione all'ordine nelle serie!]
Il Teorema di Riemann sulle serie dice che se una serie converge, ma non assolutamente (es. $\sum (-1)^n / n$), riordinando opportunamente i termini puoi farla convergere a \emph{qualsiasi} numero reale tu voglia! Per questo l'assoluta convergenza (o il segno costante) è cruciale per poter maneggiare le serie come se fossero somme normali.
\end{hint}

\chapter*{Bibliografia}
\addcontentsline{toc}{chapter}{Bibliografia}
\begin{itemize}
    \item [1] D. Acemoglu, \emph{Introduction to Modern Economic Growth}, Princeton University Press, 2009.
    \item [2] Q. Berger, F. Caravenna, and P. Dai Pra, \emph{Probabilità}, Springer, 2021.
    \item [3] D. Bertsekas, \emph{Dynamic Programming and Optimal Control}, Athena Scientifica, 2005.
    \item [5] P. Billingsley, \emph{Probability and Measure}, Wiley, 1995.
    \item [6] T. Bjork, \emph{Arbitrage theory in continuous time}, Oxford University Press, 2009.
    \item [7] P. Brémaud, \emph{Markov Chains}, Springer, 1999.
    \item [13] A. Pascucci and W. Runggaldier, \emph{Finanza Matematica: Teoria e Problemi per Modelli Multi-periodali}, Springer, 2009.
\end{itemize}


\end{document}